\documentclass[11pt]{article}

% Shared preamble for the SKJS soundness corpus (main paper + deep dives).
% Loaded via \input in each document's preamble. Documents may define
% additional packages, theorem environments, and macros AFTER this input,
% but must not redefine anything here.
\usepackage{amsmath}
\usepackage{amssymb}
\usepackage{amsthm}
\usepackage{mathtools}
\usepackage{stmaryrd}
\usepackage[margin=1in]{geometry}
\usepackage{tcolorbox}
\usepackage{booktabs}
\usepackage{array}
\usepackage{longtable}
\usepackage[hyphens]{url}
\urlstyle{tt}

% Long typewriter tokens (paths, identifiers) defeat hyphenation;
% allow inter-word stretch instead of overfull lines.
\emergencystretch=2em

\newtheorem{theorem}{Theorem}
\newtheorem{lemma}[theorem]{Lemma}
\newtheorem{definition}[theorem]{Definition}
\newtheorem{corollary}[theorem]{Corollary}
\newtheorem{axiom}[theorem]{Axiom}
\newtheorem{obligation}[theorem]{Obligation}

% Summary box for key definitions
\newtcolorbox{keydef}{
  colback=gray!5,
  colframe=gray!50,
  boxrule=0.5pt,
  arc=2pt,
  left=4pt,right=4pt,top=2pt,bottom=2pt
}

% Core semantic notation shared by every document
\newcommand{\sem}[1]{\llbracket #1 \rrbracket}
\newcommand{\Val}{\mathsf{Val}}

% Implementation anchor: file / class / function at the pinned revision
\newcommand{\impl}[1]{\par\noindent{\footnotesize\textsf{impl:} \texttt{#1}}}

% Single source for the modeled revision. Every document's \date uses this;
% the ledger and tooling read the same values.
\newcommand{\skjspinhash}{9108f339}
\newcommand{\skjspindate}{2026-07-02}
\newcommand{\skjspin}{Models revision \texttt{\skjspinhash} of \texttt{main} (\skjspindate)}

\usepackage{mathpartir}

% Core semantic notation (main-paper specific)
\newcommand{\Pred}{\mathcal{P}(\Val)}
\newcommand{\refineop}{\mathsf{refine}}

% Type domains
\newcommand{\ObjVal}{\mathsf{ObjVal}}
\newcommand{\String}{\mathsf{String}}
\newcommand{\SymVal}{\mathsf{SymVal}}
\newcommand{\SType}{\mathsf{SType}}
\newcommand{\Tags}{\mathsf{Tags}}
\newcommand{\TType}{\mathsf{TType}}
\newcommand{\Clauses}{\mathsf{Clauses}}
\newcommand{\Clause}{\mathsf{Clause}}
\newcommand{\STypes}{\mathsf{STypes}}
\newcommand{\TInst}{\mathsf{TInst}}
\newcommand{\TInsts}{\mathsf{TInsts}}
\newcommand{\TParam}{\mathsf{TParam}}
\newcommand{\TFields}{\mathsf{TFields}}
\newcommand{\TField}{\mathsf{TField}}

% Tag functions
\newcommand{\tagfn}{\mathsf{tag}}
\newcommand{\tags}{\mathsf{tags}}
\newcommand{\rtag}{\mathsf{rtag}}

% Values
\newcommand{\Payload}{\mathsf{Payload}}

% Substitution
\newcommand{\Subst}{\mathsf{Subst}}
\newcommand{\SVOpen}{\mathsf{SVOpen}}
\newcommand{\SVClosed}{\mathsf{SVClosed}}

% Environments and refinement
\newcommand{\Env}{\Gamma}
\newcommand{\RefineT}[3]{#1 \vdash #2 \stackrel{\mathsf{t}}{\rightsquigarrow} #3}
\newcommand{\RefineF}[3]{#1 \vdash #2 \stackrel{\mathsf{f}}{\rightsquigarrow} #3}
\newcommand{\RefineB}[4]{#1 \vdash #2 \stackrel{#4}{\rightsquigarrow} #3}

% Clause notation
\newcommand{\clause}[2]{\langle #1 \mid #2 \rangle}

\title{Refinement and Inference Growth in SKJS:\\ a Soundness Model of the Type Checker}
\author{}
\date{\skjspin}

\begin{document}

\maketitle

\begin{abstract}
We formalize the internal type representation used by SKJS and prove soundness
properties for its two flow-sensitive type transformations: predicate-driven
\emph{refinement} (control-flow narrowing, \texttt{evalCond}) and
\emph{constant removal} (the residue of the former general type-difference
operation). We model the inference substitution as a small state machine and
state the invariants under which inference growth (\texttt{growTVars}) cannot
cause unsound acceptance. Everything the model does not describe concretely is
either an explicit abstract interface with stated axioms (collected in
Section~\ref{sec:assumptions}, Assumptions and Proof Obligations) or explicitly
out of scope with a justification that the omission cannot hide a soundness
hole (Section~\ref{sec:scope}). This document supersedes the draft
attached to PR \#342 (``Refinement and Type Difference in SKJS'').
\end{abstract}

\tableofcontents

%==========================================================================
\section{Scope, Method, and Modeling Contract}
\label{sec:scope}

\paragraph{Pinned revision.}
All implementation references (file, class, function) are against commit
\texttt{\skjspinhash}. The document is maintained by the discipline of
Appendix~\ref{app:maintenance}: every inference rule in the refinement master
list corresponds to one dispatch arm of \texttt{evalCond}, and drift between
the list and the dispatch is intended to be mechanically checkable.
Two terms of art for that mechanism, used throughout: an
\emph{implementation anchor} is a file/class/function reference (typeset as
\textsf{impl:}~\texttt{\dots}) locating the code a definition or rule
models — the maintenance check verifies anchors bidirectionally; a
\emph{pinned test} is a test in the repository that fails when a documented
behavior changes, forcing the corresponding passage to be updated in the
same change.

\paragraph{Modeling contract.}
Every behavior of the implementation that can affect type soundness is placed
in exactly one of three bins:
\begin{enumerate}
\item \textbf{Modeled}: defined concretely in this document and proved sound.
\item \textbf{Abstracted}: hidden behind an interface with explicitly stated
  axioms. All such axioms are collected in Section~\ref{sec:assumptions}
  (Assumptions and Proof Obligations), each with an implementation anchor
  and a status from the fixed vocabulary \textsc{proved} / \textsc{assumed}
  / \textsc{open} / \textsc{violated} defined there.
\item \textbf{Out of scope}: listed below, each with a reason why omission
  cannot create a hole in the soundness argument.
\end{enumerate}
An implementation behavior that fits none of the bins is a defect of this
document.

\paragraph{Terminology.}
The model deliberately keeps the implementation's names so that every term
can be grepped for; the cost is a few SKJS-specific usages, fixed here.
An \emph{object branch} is one alternative of an object type's internal
union (the implementation's \texttt{TObject.branches}); it is always
qualified, since \emph{true/false branch} refers to control flow.
A \emph{member} is one $\SType$ alternative of an $\STypes$ union;
a \emph{clause} pairs type parameters with such a union.
A \emph{candidate} is an ordinary inference variable (\texttt{TVID});
\emph{grow} and \emph{seal} are the substitution transitions of
Section~\ref{sec:inference}; \emph{expansion} rewrites a type to its
instance-plus-clauses view under the current substitution.
A \emph{brand} is the \texttt{Array.isArray} marker on an object branch;
a \emph{phantom field} records runtime-confirmed presence without a
declaration. ``Open'' is used only for the substitution state
$\SVOpen$ and the obligation status; the object semantics of
Definition~\ref{def:object-semantics} is called \emph{extensible}.

\paragraph{Out of scope, with justification.}
\begin{itemize}
\item \emph{Syntax rejected at normalization.} Conditional types
  (\texttt{T extends U ? X : Y}, SKJS9020), \texttt{infer} (9017), template
  literal types (9018), indexed access \texttt{T[K]} (9021), type predicates
  \texttt{x is T} (9015), \texttt{asserts} predicates (9013): these produce a
  fatal diagnostic in \texttt{TEnv::normalize} and never reach the type
  machinery, so they are outside the soundness surface by construction.
  \impl{TEnv.sk / normalize}
\item \emph{Unhandled condition forms.} The refinement dispatch falls through
  to the identity on any condition shape it does not recognize
  (Rule \textsc{R-Other}), and the identity refinement is trivially sound
  (Lemma~\ref{lem:identity-sound}). Consequently a stale or incomplete master
  list makes this document \emph{incomplete}, never \emph{unsound}, as a
  description of narrowing.
\item \emph{Mutation during checking of straight-line code.} This document
  proves \emph{one-step branch soundness}: a single condition is evaluated
  against a fixed store. Assignment, loops, and join points are treated only
  through the union-soundness of scope merging
  (Lemma~\ref{lem:merge-sound}); a full small-step preservation theorem is
  future work. The interaction of narrowing with mutation \emph{during
  condition evaluation} is governed by assumptions
  \textbf{L-PURE} and \textbf{L-LOCAL}.
\item \emph{Flow-typing sites outside condition narrowing.} The master list
  covers condition forms only. The checker performs flow-sensitive type
  updates at five further sites, listed here so the boundary is a stated
  enumeration rather than an implication: assignment to a local updates its
  scope entry (\texttt{TypeCheckExpr.sk}, \texttt{Var}-LHS
  \texttt{scope.set}); literal-typed local objects grow fields on
  assignment; \texttt{this} is re-typed after \texttt{super(...)}; the
  non-null assertion \texttt{x!} and optional chaining remove nullish
  members from the \emph{expression's} type (no scope mutation, via
  \texttt{removeUndefined}); and the \texttt{for}-\texttt{in} index
  mechanism suppresses element errors without narrowing. None of these is a
  condition refinement; each would be a case of the future preservation
  theorem. Note also that \texttt{for}-loop conditions do \emph{not} narrow.
\item \emph{Numeric coercions, operators, and standard-library semantics.}
  Only the condition forms of the master list are given evaluation semantics.
\end{itemize}

%==========================================================================
\section{Runtime Model}
\label{sec:runtime}

\begin{definition}[Runtime Tags]
\label{def:runtime-tags}
$\Tags$ is the nine-element set
\[
\Tags = \{\mathsf{undef}, \mathsf{void}, \mathsf{null}, \mathsf{bool},
          \mathsf{num}, \mathsf{str}, \mathsf{bigint}, \mathsf{sym},
          \mathsf{obj}\}.
\]
\impl{TType.sk / STypeTag (TagUndefined \dots\ TagObject)}
\end{definition}

There is no dedicated tag for tuples, promises, or algebraic data types:
tuples and promises are object-tagged values whose types live in the
instance layer or as object types
(\texttt{SKJS\_Tuple}/\texttt{ReadonlyTuple} instances; \texttt{Promise} as a
stub-declared object type), and discriminated unions are ordinary object
types with discriminant fields.

\begin{definition}[Runtime Values]
\label{def:runtime-values}
Fix payload sets $(\Payload_k)_{k\in\Tags}$ and define
$\Val \triangleq \{(k,p) \mid k \in \Tags,\ p \in \Payload_k\}$, with
$\rtag((k,p)) = k$. Primitive payloads are concrete; object payloads are
abstract:
\begin{align*}
\Payload_{\mathsf{undef}} &= \{\star\} &
\Payload_{\mathsf{void}} &= \{\star\} &
\Payload_{\mathsf{null}} &= \{\star\} \\
\Payload_{\mathsf{bool}} &= \{\mathsf{true},\mathsf{false}\} &
\Payload_{\mathsf{num}} &= \mathbb{F} &
\Payload_{\mathsf{str}} &= \String \\
\Payload_{\mathsf{bigint}} &= \mathbb{Z} &
\Payload_{\mathsf{sym}} &= \SymVal &
\Payload_{\mathsf{obj}} &= \Payload'_{\mathsf{obj}}
\end{align*}
where $\mathbb{F}$ is the set of IEEE-754 doubles (including $-0$ and
$\mathsf{NaN}$), $\SymVal$ an abstract infinite set of symbol identities, and
$\Payload'_{\mathsf{obj}}$ an arbitrary set. The tag $\mathsf{void}$ models
the checker's distinct \texttt{void} type; at runtime its unique payload is
the value \texttt{undefined}, which is why refinements treat the tags
$\mathsf{undef}$ and $\mathsf{void}$ together (Definition~\ref{def:eval}).
\end{definition}

\begin{definition}[Object Interface]
\label{def:object-interface}
Let $\ObjVal \triangleq \{v \in \Val \mid \rtag(v) = \mathsf{obj}\}$.
We assume the following abstract interface on object values:
\begin{itemize}
\item Keys: $\mathsf{Key} \triangleq \String \uplus \SymVal$.
\item $\mathsf{props} : \ObjVal \to \mathcal{P}(\mathsf{Key})$ and a partial
  $\mathsf{get} : \ObjVal \times \mathsf{Key} \rightharpoonup \Val$, with
  $\mathsf{get}(o,p)$ defined iff $p \in \mathsf{props}(o)$. Write
  $\mathsf{has}(o,p)$ for $p \in \mathsf{props}(o)$.
\item $\mathsf{callable}, \mathsf{constructible} : \ObjVal \to
  \{\mathsf{true},\mathsf{false}\}$, with the JS fact
  $\mathsf{constructible}(o) \implies \mathsf{callable}(o)$ (every
  constructor is a function; \texttt{typeof} yields \texttt{"function"} for
  it).
\item $\mathsf{isArrayV} : \ObjVal \to \{\mathsf{true},\mathsf{false}\}$, the
  runtime truth of \texttt{Array.isArray}.
\item A classification $\mathsf{kind} : \String \to
  \{\mathsf{strKey}, \mathsf{numKey}\}$ of string keys into numeric and
  non-numeric spellings, used by index signatures.
\end{itemize}
This interface is a snapshot view: $\mathsf{get}$ describes the properties of
$o$ at the moment a condition is evaluated. See assumption \textbf{L-PURE}.
\end{definition}

\begin{definition}[Strict Equality]
\label{def:seq}
JS strict equality is \emph{not} identity on $\Val$; we model it as a
relation ${\approx} \subseteq \Val \times \Val$:
$v \approx w$ iff one of
\begin{itemize}
\item $\rtag(v), \rtag(w) \in \{\mathsf{undef}, \mathsf{void}\}$ (the
  \texttt{void} tag is a checker fiction for the runtime value
  \texttt{undefined});
\item $\rtag(v) = \rtag(w) = \mathsf{num}$ and the payloads are equal as
  IEEE doubles with $-0$ identified with $0$ and $\mathsf{NaN}$ equal to
  nothing (including itself);
\item $\rtag(v) = \rtag(w) \notin \{\mathsf{undef},\mathsf{void},
  \mathsf{num}\}$ and the payloads are equal (reference identity for
  objects).
\end{itemize}
Note $\approx$ is not reflexive ($\mathsf{NaN} \not\approx \mathsf{NaN}$)
and crosses the $\mathsf{undef}/\mathsf{void}$ tag pair. A set
$X \subseteq \Val$ is \emph{$\approx$-closed} if $v \in X$ and $w \approx v$
imply $w \in X$.
\end{definition}

\begin{lemma}[$\approx$ is a Partial Equivalence]
\label{lem:per}
$\approx$ is symmetric and transitive; its domain (the values related to
themselves) is $\Val$ minus the $\mathsf{NaN}$-payload numbers. On that
domain ``the $\approx$-class of $v$'' is well-defined.
\end{lemma}
\begin{proof}
Each arm is symmetric. Transitivity: a chain stays within the
$\{\mathsf{undef},\mathsf{void}\}$ pair or within a single tag with
payload equality (for numbers, equality-with-$\pm 0$-identified is
transitive and $\mathsf{NaN}$ relates to nothing, so chains through it are
vacuous).
\end{proof}

\begin{definition}[Truthiness]
\label{def:truthiness}
$\mathsf{truthy} : \Val \to \{\mathsf{true},\mathsf{false}\}$ is defined
per tag: values of tags $\mathsf{undef}$, $\mathsf{void}$, $\mathsf{null}$
are falsy; $(\mathsf{bool}, b)$ is truthy iff $b = \mathsf{true}$;
$(\mathsf{num}, n)$ is falsy iff $n \in \{0, -0, \mathsf{NaN}\}$;
$(\mathsf{str}, s)$ is falsy iff $s = \texttt{""}$;
$(\mathsf{bigint}, z)$ is falsy iff $z = 0$;
symbols and objects are always truthy.
\impl{TType.sk / TPrimValueKind::isFalsy, SType::falsiness}
\end{definition}

%==========================================================================
\section{Internal Type Representation}
\label{sec:representation}

\subsection{Simple types}

\begin{definition}[Simple Types]
\label{def:stype}
$\SType$ has exactly three constructor families:
\begin{itemize}
\item $\mathsf{unit}(u)$ for $u \in \{\texttt{undefined}, \texttt{void},
  \texttt{null}\}$;
\item $\mathsf{prim}(k, L)$ for
  $k \in \{\texttt{bool}, \texttt{num}, \texttt{str}, \texttt{bigint},
  \texttt{sym}\}$ and $L \in \{\top\} \cup
  \mathcal{P}^{+}_{\mathrm{fin}}(\mathsf{Lit}_k)$, where $\mathsf{Lit}_k$ is
  the set of literal spellings of kind $k$ and $\top$ means the full
  primitive;
\item $\mathsf{object}(B)$ for a nonempty finite set $B$ of \emph{branches}
  (Definition~\ref{def:fields}).
\end{itemize}
$\tagfn$ maps each constructor to the evident tag.
\impl{TType.sk / TPrimUnit, TPrimValue, TObject; SType::tag}
\end{definition}

\begin{definition}[Literal Denotation]
\label{def:literal-denotation}
Assume a denotation $\delta_k : \mathsf{Lit}_k \to \Payload_{\tagfn(k)}$ from
literal spellings to payloads, and define
\begin{align*}
\sem{\mathsf{unit}(u)} &\triangleq \{(\tagfn(u), \star)\} \\
\sem{\mathsf{prim}(k, \top)} &\triangleq
  \{(\tagfn(k), p) \mid p \in \Payload_{\tagfn(k)}\} \\
\sem{\mathsf{prim}(k, L)} &\triangleq
  \{\, v \mid \rtag(v) = \tagfn(k),\
     \exists \ell \in L.\ v \approx (\tagfn(k), \delta_k(\ell)) \,\}.
\end{align*}
Literal-type denotations are the \emph{same-tag $\approx$-closure} of their
spellings' denotations: the literal type \texttt{0} contains both $0$ and
$-0$, matching JS strict equality and tsc's treatment of numeric literal
types. (For non-numeric kinds the closure is the identity. The cross-tag
$\mathsf{undef}/\mathsf{void}$ arm of $\approx$ is deliberately \emph{not}
baked into unit denotations; the rules that compare against
\texttt{undefined} handle the two tags jointly, and the residual corner is
assumption \textbf{L-SEQ}.)
We assume $\mathsf{Lit}_{\texttt{num}}$ contains no $\mathsf{NaN}$
spelling: there is no NaN literal type, and a hypothetical one would denote
$\emptyset$ under closure semantics (nothing is ${\approx}\,\mathsf{NaN}$),
making \textbf{L-TC} unsatisfiable for it.
The implementation stores literal \emph{spellings} and compares them
syntactically; assumption \textbf{L-LIT} states the assumptions this
requires. Unique symbols are represented as literals of kind \texttt{sym}
whose stored string is a \emph{synthesized position-derived identity}
(\texttt{unique\_symbol:file:sym:offset:name}), not a source lexeme — so
distinct unique symbols have distinct denotations by construction.
\impl{PrimLiterals.sk; TType.sk / TPrimValue; TEnv.sk / uniqueSymbol}
\end{definition}

\begin{lemma}[A0: Tag Consistency]
\label{lem:a0}
$\forall s \in \SType,\ v \in \sem{s}.\ \rtag(v) = \tagfn(s)$.
\end{lemma}
\begin{proof}
Immediate from the definitions above and
Definition~\ref{def:object-semantics}, whose members are drawn from
$\ObjVal$.
\end{proof}

\begin{lemma}[A1: Tag Disjointness]
\label{lem:a1}
$\forall s_1, s_2 \in \SType.\
\tagfn(s_1) \neq \tagfn(s_2) \implies \sem{s_1} \cap \sem{s_2} = \emptyset$.
\end{lemma}
\begin{proof}
By A0, as values carry a unique tag.
\end{proof}

\subsection{Object branches and fields}

\begin{definition}[Fields and Branches]
\label{def:fields}
A \emph{field} is a record
$f = (\mathsf{name}, \mathsf{opt}, \mathsf{get}, \mathsf{set},
\mathsf{access})$
with $\mathsf{name} \in \mathsf{Key}$, $\mathsf{opt} \in
\{\mathsf{true},\mathsf{false}\}$, $\mathsf{get}, \mathsf{set} \in \TType$,
and $\mathsf{access} \in \{\mathsf{Pub}, \mathsf{Priv}(\mathit{owner}),
\mathsf{Never}, \mathsf{Phantom}\}$.
An \emph{object branch} $b \in \TFields$ is a record consisting of:
\begin{itemize}
\item finite maps $\mathsf{flds}, \mathsf{symflds}$ from string/symbol keys
  to fields;
\item optional index fields $\mathsf{idx}_\kappa$ for
  $\kappa \in \{\mathsf{strKey}, \mathsf{numKey}, \mathsf{symKey}\}$, each a
  field;
\item $\mathsf{call}, \mathsf{new} \in
  \{\mathsf{None}\} \cup \{(\mathit{opt}, \sigma) \mid \mathit{opt} \in
  \{\mathsf{true},\mathsf{false}\}\}$ where $\sigma$ ranges over abstract
  function signatures;
\item $\mathsf{isArr} \in \{\mathsf{None}, \mathsf{Some}(\mathsf{true}),
  \mathsf{Some}(\mathsf{false})\}$, the array brand.
\end{itemize}
Readonly fields are encoded as $\mathsf{set} = \bot$
(\texttt{TField::create} with \texttt{isReadonly}).
The implementation's \texttt{TObject} additionally carries
\texttt{unexpandedType :\ ?TSInst}, a best-effort record of the named type
an object came from; it is a printing hint plus one program-\emph{restricting}
check (the async-return ``must be global \texttt{Promise}'' rule) and never
participates in $\sem{\cdot}$, subtyping, or refinement, so the model omits
it.
\impl{TType.sk / TField, TFields, TIndexesBase, FieldAccess}
\end{definition}

\begin{definition}[Signature Modes]
\label{def:sig-modes}
The three-valued call/new encoding has the following intended reading
(abbreviate: $\mathsf{None}$ / $\mathsf{Maybe}(\sigma)$ /
$\mathsf{Definite}(\sigma)$):
$\mathsf{None}$ = definitely not callable (constructible);
$(\mathsf{true}, \sigma)$ = \emph{maybe} callable, with signature $\sigma$ if
so; $(\mathsf{false}, \sigma)$ = definitely callable with signature $\sigma$.
The implementation violates the $\mathsf{None}$ reading at
class-constructor \emph{values}: they are typed $\mathsf{None}$ although
class constructors implement $[[\mathrm{Call}]]$
(\texttt{typeof C === "function"} at runtime), so
\textsc{R-Typeof}'s soundness proof — which consumes this reading —
inherits a witnessed violation at those values (ledger row A-CTOR-CALL;
issue draft \path{class-constructor-call-none-typeof.md}; deep dive A).
The same $\mathsf{None}$ mode is read as ``callability unknown'' by one
subtyping arm (deep dive D's F-CALL) — the \texttt{call{=}None} dichotomy:
no single reading validates both consumers.
\end{definition}

\begin{definition}[Object-Type Semantics]
\label{def:object-semantics}
$o \in \sem{\mathsf{object}(B)}$ iff $o \in \ObjVal$ and there exists an
object branch $b \in B$ with $o \in \sem{b}$, where $o \in \sem{b}$ iff all
of:
\begin{itemize}
\item \textbf{Required fields:} for every $f \in
  \mathsf{flds}(b) \cup \mathsf{symflds}(b)$ with
  $\mathsf{opt}(f) = \mathsf{false}$: $\mathsf{has}(o, \mathsf{name}(f))$ and,
  unless $\mathsf{access}(f) = \mathsf{Phantom}$,
  $\mathsf{get}(o, \mathsf{name}(f)) \in \sem{\mathsf{get}(f)}$.
\item \textbf{Optional fields:} for every such $f$ with
  $\mathsf{opt}(f) = \mathsf{true}$: if $\mathsf{has}(o, \mathsf{name}(f))$
  then $\mathsf{get}(o, \mathsf{name}(f)) \in \sem{\mathsf{get}(f)}$.
\item \textbf{Indexes:} an index constrains the props it \emph{covers}:
  the string index covers every string key (of either $\mathsf{strKey}$ or
  $\mathsf{numKey}$ kind), the number index covers $\mathsf{numKey}$-kind
  keys only, and the symbol index covers symbol keys — matching TS, where a
  string index signature also constrains numeric properties, and the
  implementation's numeric lookup (\texttt{numberIndex}~\texttt{??}
  \texttt{stringIndex}). For each index
  $\mathsf{idx}_\kappa(b) = \mathsf{Some}(f)$ and every
  $p \in \mathsf{props}(o)$ covered by it and not named by a field of $b$:
  $\mathsf{get}(o, p) \in \sem{\mathsf{get}(f)}$.
\item \textbf{Call/new:} if $\mathsf{call}(b) = \mathsf{None}$ then
  $\neg\mathsf{callable}(o)$; if $\mathsf{call}(b) = (\mathsf{false},\sigma)$
  then $\mathsf{callable}(o)$; if $\mathsf{call}(b) = (\mathsf{true},\sigma)$,
  no constraint. Signatures themselves are opaque (their behavioral meaning is
  bundled into assumption \textbf{L-SUB}). Same for $\mathsf{new}$ with
  $\mathsf{constructible}$.
\item \textbf{Brand:} if $\mathsf{isArr}(b) = \mathsf{Some}(x)$ then
  $\mathsf{isArrayV}(o) = x$.
\end{itemize}
The semantics is \emph{extensible}: properties not mentioned are
unconstrained. A
$\mathsf{Phantom}$ field constrains only \emph{presence}, not the value — it
records knowledge obtained from a runtime \texttt{in} check when no
declaration is known; the implementation additionally makes phantoms error at
access and be invisible to subtyping, which only \emph{restricts} programs and
so needs no semantic counterpart. The same restrict-only argument covers
$\mathsf{Never}$ (produced when two same-named private fields with different
owners meet: accessing it is a static error, but as a membership constraint
it is an ordinary field). Note that membership constrains values via
$\mathsf{get}$-types only; the $\mathsf{set}$ component governs which writes
the checker accepts and participates in soundness through subtyping and
mutation, not through membership (assumption \textbf{L-MUT}). That
governance has a second named facet, cross-view coherence: a checked
field or index write through one mutable view preserves the membership
guarantees of every live alias of the same runtime object
(\textbf{L-MUT-ALIAS}; its call-mediated counterpart is the row
L-MUT-ALIAS-CALL) — which holds only while every mutable view's
$\mathsf{set}$-types stay at or below the object's actual write
capability.
For recursive declarations the membership conditions recurse — a field's
$\sem{\mathsf{get}(f)}$ can mention $\sem{b}$ again — and every recursive
occurrence is positive, so least and greatest solutions exist
(Knaster--Tarski); $\sem{\cdot}$ is the \emph{greatest} solution, the
reading under which a cyclic value (\texttt{const x = \{\}; x.next = x})
inhabits its recursive type. This choice is registry row \textbf{SEM-GFP}.
\impl{TType.sk / TObject (interpretation comment), TFields}
\end{definition}

\begin{definition}[Write Authorization]
\label{def:write-auth}
A view $b$ (an object branch, as recorded for a local or a field)
\emph{authorizes} the write of value $v$ at key $p$ iff
$v \in \sem{\mathsf{setty}(b,p)}$, where $\mathsf{setty}(b,p)$ is the
resolved $\mathsf{set}$-type of $p$ in $b$: the intersection
$\bigsqcap f.\mathsf{set}$ of a named field $f$, else the covering
index's $\mathsf{set}$-type resolved as at lvalue sites (the number
index falling back to the string index); $\mathsf{setty}(b,p) = \bot$
when $p$ has no writable field or covering index, and for a readonly
field ($\mathsf{set} = \bot$) — either way $b$ authorizes nothing at
$p$. Optionality does not weaken authorization: an accepted write makes
$p$ present, and the stored value must satisfy the field's get-clause of
Definition~\ref{def:object-semantics}. The checker realizes this
definition at assignment
sites — the right-hand side is checked against $\mathsf{setty}$ by
subtyping (\textbf{L-SUB}) — and writes mediated by method calls are
covered by the signature semantics (E-SIG), not by $\mathsf{setty}$.
This is the corpus's definition of a view's write surface, registry row
\textbf{SET-SEM}; \textbf{L-MUT-ALIAS} is the coherence property over
it.
\impl{TType.sk / TFields::resolveFieldType (isLValue);
TypeCheckExpr.sk / assign}
\end{definition}

\subsection{The union spine}

\begin{definition}[STypes]
\label{def:stypes}
$\STypes$ is a finite, tag-indexed set: at most one $\SType$ per tag;
$\sem{S} = \bigcup_{s \in S}\sem{s}$. Write $S[k]$ for the representative at
tag $k$ if present.
\impl{TType.sk / STypes (InternalKeyedArrayMap keyed by STypeTag)}
\end{definition}

Precision within a tag is carried by payloads: literal sets for primitives
and object-branch sets for objects. A union of object types is represented
exactly — one object branch per alternative — so merging alternatives loses
no precision at the object tag.

\begin{definition}[Instances]
\label{def:tinsts}
$\TInst$ is the abstract instance layer: type variables $\alpha$
(\texttt{TVar}, carrying a set of constraint types), unexpanded named
instantiations (\texttt{TSInst}), unexpanded intersections (\texttt{TInter}),
and delayed array-literal forms (\texttt{TArrayLit}). $\TInsts$ is a finite
set of instances with $\sem{I}_\rho = \bigcup_{i \in I}\sem{i}_\rho$. Instance
semantics is \emph{parametrized by a substitution state} $\rho$
(Section~\ref{sec:inference}).  The denotation of a type-variable occurrence
$\alpha^{cs}$ (identifier $\alpha$ carrying the occurrence-level constraint
set $cs$) is
\[
  \sem{\alpha^{cs}}_\rho = \mathcal E_\rho(\alpha) \cap
    \bigcap_{c\in cs}\sem{c}_\rho,
  \qquad
  \mathcal E_\rho(\alpha) =
  \begin{cases}
    \emptyset & \alpha\notin\operatorname{dom}(\rho),\\
    \sem{t}_\rho & \rho(\alpha)\in
      \{\SVOpen(t),\SVClosed(t)\}.
  \end{cases}
\]
A binding type $t$ may itself mention type variables, so this equation
joins the recursive membership conditions of
Definition~\ref{def:object-semantics}; every recursive occurrence is
positive, and the solution is the greatest one (registry row
SEM-GFP). An unbound \texttt{TVar} thus denotes an empty inference lower
bound, and \texttt{TVar} is the model's only variable kind — the model
has no rigid or skolem variables. Constraints belong to
occurrences because \texttt{TVar::setConstraint} is functional;
\texttt{Subst::get} realizes the displayed intersection when the
occurrence is read. This equation is the corpus's definition of
\texttt{TVar} denotation, registry row \textbf{TVAR-SEM}. Named
instances and intersections satisfy the expansion axiom \textbf{L-EXP}.
\impl{TType.sk / TInst, TVar, TSInst, TArrayLit; TTypeInter.sk / TInter}
\end{definition}

\begin{definition}[Clause, Clauses, TType]
\label{def:clauses}
\leavevmode
\begin{keydef}
A clause $\clause{V}{S}$ pairs type parameters $V \subseteq_{\mathrm{fin}}
\TParam$ with $S \in \STypes \cup \{\top\}$; clauses cannot be bottom
(an empty $\STypes$ is dropped at construction). $\Clauses$ is a map from
var-sets to clauses, at most one per key, with the invariant that if the top
clause $\clause{\emptyset}{\top}$ is present it is alone. $\TType = (I, C)$
with $I \in \TInsts$, $C \in \Clauses$.
\[
\sem{\clause{V}{S}}_\rho = \bigcap_{v \in V}\sem{v}_\rho \cap \sem{S}_\rho
\qquad
\sem{C}_\rho = \bigcup_{c \in C}\sem{c}_\rho
\qquad
\sem{(I,C)}_\rho = \sem{I}_\rho \cup \sem{C}_\rho
\]
where $\sem{\top}_\rho = \Val$.
\end{keydef}
A clause may have $V \neq \emptyset$ with
$S = \top$ (a type bounded only by its parameters); only
$\clause{\emptyset}{\top}$ is the top clause.
\impl{TType.sk / Clause, Clauses, TType}
\end{definition}

\subsection{Operation interfaces}
\label{sec:op-interfaces}

Refinement and constant removal are built from a small set of operations whose
soundness this document assumes as assumptions and uses freely:

\begin{axiom}[L-UNION, L-INTER]
\label{ax:union-inter}
For all $t_1, t_2$ and states $\rho$:
$\sem{t_1}_\rho \cup \sem{t_2}_\rho \subseteq
 \sem{\mathsf{union}(t_1,t_2)}_\rho$
and
$\sem{t_1}_\rho \cap \sem{t_2}_\rho \subseteq
 \sem{\mathsf{inter}(t_1,t_2)}_\rho$.
In particular $\mathsf{inter}(t_1,t_2) = \bot$ implies
$\sem{t_1}_\rho \cap \sem{t_2}_\rho = \emptyset$, which refinement uses to
drop branches. (The two clauses are tracked as separate registry rows:
the containment as L-INTER, the disjointness clause as L-INTER-BOT.) The same containments are assumed at the $\STypes$ and
$\mathsf{object}$ levels.
\impl{TTypeUnion.sk; TTypeInter.sk; TType.sk / STypes::union, TObject::inter}
\end{axiom}

\begin{axiom}[L-EXP: Expansion]
\label{ax:expansion}
Expansion rewrites a type to an instance-plus-clauses view using the current
substitution state, possibly evolving that state
($\rho \rightarrow \rho'$, e.g.\ sealing a fresh covariantly-read variable to
top). We assume
$\sem{t}_\rho \subseteq \sem{\mathsf{expand}(t)}_{\rho'}$ and that residual
instances (those left in the instance layer after expansion) keep their
semantics. All refinement operators run on the expanded view. Expansion
takes a mode (\texttt{expandTVars}); the containment is claimed for modes
\texttt{None}/\texttt{Covariant} only — contravariant expansion is an
under-approximating read belonging to \textbf{L-SUB} — and every
refinement client uses those modes (deep dive C).
\impl{TType.sk / TType::expand, expandToVectors, mapSTypes}
\end{axiom}

\begin{axiom}[L-MONO: Refinement-Time State Monotonicity]
\label{ax:mono}
The substitution evolutions performed \emph{while refining a condition}
(expansion reads, which may seal fresh variables to top, and exploratory
intersections) preserve membership for the types \emph{installed in the
environment} at the evolution, and for the type the rule just produced: if
$\rho \to \rho'$ is such an evolution, then for every
$u \in \mathrm{range}(\Env)$,
$v \in \sem{u}_\rho \implies v \in \sem{u}_{\rho'}$.
Without this, narrowing one variable could silently invalidate the recorded
types of \emph{other} variables. Two remarks. First, deep dive E proves
this axiom outright (given semantic monotonicity in variable envelopes and
\textbf{L-UNION}); in particular there is no constraint-mutation threat
from shared state — \texttt{TVar::setConstraint} is a \emph{functional}
update on a frozen value, so constraints are occurrence-level (registry
row SETCONSTRAINT-FUN) — and the
$\mathrm{range}(\Env)$ restriction, stated here for safety, is
droppable. Second, the classic
contravariance counterexample (function parameter positions) is deflected
only because signatures are opaque to $\sem{\cdot}$ — membership never
inspects parameter types (Definition~\ref{def:object-semantics}). This is
the refinement-time counterpart of Obligation~\ref{obl:monotone} (which
concerns inference-time growth), and it is used by every branch-soundness
argument below.
\impl{TTypeSubst.sk / Subst::get (seal-to-top); TypeCheckRefine.sk
(cloned-subst intersections)}
\end{axiom}

%==========================================================================
\section{Refinement (Control-Flow Narrowing)}
\label{sec:refinement}

\subsection{The single refinement primitive}

Nearly all narrowing in SKJS is performed by one primitive applied to a
local variable's type; the exceptions (equality's whole-variable
intersection and the connectives' scope merge) are called out after the
framework lemmas and receive direct proofs. The primitive is a filter/map:
a partial mapper transforms each $\SType$ member of each clause of the
expanded type, and an optional default replaces each \emph{top} clause.

\begin{definition}[Memberwise Refinement]
\label{def:branchwise}
Let $g : \SType \rightharpoonup \SType$ be a partial, tag-preserving mapper
and $d \in \{\mathsf{None}\} \cup \{\mathsf{Some}(D) \mid D \in \STypes\}$ an
optional default. Define $\refineop^{g,d}(t)$ as follows: expand $t$
(axiom \textbf{L-EXP}) to residual instances $I$ and clauses $C$; replace each
clause $\clause{V}{S}$ by
\[
\begin{cases}
\clause{V}{D} & \text{if } S = \top \text{ and } d = \mathsf{Some}(D)\\
\clause{V}{\top} & \text{if } S = \top \text{ and } d = \mathsf{None}\\
\clause{V}{\{g(s) \mid s \in S,\ g(s)\ \mathrm{defined}\}} & \text{otherwise},
\end{cases}
\]
dropping clauses whose $\STypes$ becomes empty; the result is
$(I, C')$ rebuilt with the clause union.
\impl{TypeCheckRefine.sk / setLocalFilterMap; TType.sk / TType::mapSTypes}
\end{definition}

\begin{definition}[Predicate Soundness]
\label{def:pred-sound}
Let $P \in \Pred$. A partial mapper $g$ is \emph{sound for $P$} if for every
$s \in \SType$ and $v \in \sem{s} \cap P$: $g(s)$ is defined and
$v \in \sem{g(s)}$. A default $d = \mathsf{Some}(D)$ is \emph{adequate for
$P$} if $P \subseteq \sem{D}$; $d = \mathsf{None}$ is always adequate.
Object-typed members make these denotations state-indexed (their field types
may mention instances); throughout, $\sem{s}, \sem{g(s)}, \sem{D}$ are read
at the post-refinement state $\rho'$.
A transformer $R : \TType \to \TType$ is \emph{predicate-sound for $P$} if
$\sem{t}_\rho \cap P \subseteq \sem{R(t)}_{\rho'}$ for all $t$ (with
$\rho \to \rho'$ the state evolution of $R$).
\end{definition}

\begin{lemma}[Memberwise Refinement is Predicate-Sound]
\label{lem:branchwise-sound}
If $g$ is sound for $P$ and $d$ adequate for $P$, then $\refineop^{g,d}$ is
predicate-sound for $P$.
\end{lemma}
\begin{proof}
Let $v \in \sem{t}_\rho \cap P$. By \textbf{L-EXP},
$v \in \sem{(I,C)}_{\rho'}$ for the expanded view. If $v \in \sem{I}$, the
instance layer is unchanged and $v$ survives. Otherwise
$v \in \sem{\clause{V}{S}}$ for some clause. If $S = \top$: with
$d = \mathsf{None}$ the clause is unchanged; with $d = \mathsf{Some}(D)$,
$v \in \bigcap_{V}\sem{\cdot}$ still, and $v \in P \subseteq \sem{D}$, so
$v \in \sem{\clause{V}{D}}$. If $S \neq \top$: $v \in \sem{s}$ for some
$s \in S$, and by soundness of $g$, $g(s)$ is defined with
$v \in \sem{g(s)}$, so the clause (which is then nonempty) retains $v$;
clauses dropped for emptiness never contained $v$. Tag preservation of $g$
is what makes the image a well-formed $\STypes$ (at most one member per
tag), so no two survivors merge or clobber.
Rebuilding with clause union only enlarges (\textbf{L-UNION}).
\end{proof}

\begin{lemma}[Identity is Sound]
\label{lem:identity-sound}
The identity transformer is predicate-sound for every $P$. \qed
\end{lemma}

\paragraph{Beyond the primitive.}
Two rule families do not factor through $\refineop^{g,d}$:
equality narrowing on a whole variable, which \emph{replaces} the type by an
intersection (Rule \textsc{R-Eq}); and the connective rules, which merge two
refined scopes (Rule \textsc{R-And}/\textsc{R-Or}). They get direct proofs.

\subsection{Environments, evaluation, and branch soundness}

\begin{definition}[Environments and Stores]
$\Env$ maps the \emph{local} variables of the current scope to $\TType$;
$\sigma$ maps locals to $\Val$. Write $\sigma \models_\rho \Env$ when
$\sigma(x) \in \sem{\Env(x)}_\rho$ for all $x \in \mathrm{dom}(\Env)$.
Only locals are narrowed: the implementation's refinement primitive is a
no-op for bindings not in the local scope (module-level declarations), since
narrowing a global is unsound under concurrent mutation through
closures or aliases. (Assumption \textbf{L-LOCAL}.)
\impl{TypeCheckRefine.sk / setLocalFilterMap (non-local early return)}
\end{definition}

\begin{definition}[Evaluation of Condition Forms]
\label{def:eval}
For each condition form $e$ of the master list we define relations
$\sigma \models_{\mathsf{true}} e$ and $\sigma \models_{\mathsf{false}} e$,
capturing the runtime outcomes under which the corresponding branch is taken.
Evaluation of subexpressions is assumed not to modify $\sigma$
(assumption \textbf{L-PURE}); forms whose evaluation would throw (e.g.\
\texttt{"p" in 3}) satisfy neither relation, which makes any narrowing of an
unreachable branch vacuously sound.
\begin{itemize}
\item $\sigma \models_b (e_1\ \texttt{===}\ e_2)$ iff $v_1 \approx v_2$
  (Definition~\ref{def:seq}) exactly when $b = \mathsf{true}$, where
  $v_1, v_2$ are the values of $e_1, e_2$. \texttt{!==} swaps $b$.
\item $\sigma \models_b (\texttt{typeof}\ x\ \texttt{===}\ k)$ iff
  $\sigma(x) \in P_{\mathsf{typeof}=k}$ exactly when $b = \mathsf{true}$,
  with
  \begin{align*}
  P_{\mathsf{typeof}=\texttt{"undefined"}} &=
    \{v \mid \rtag(v) \in \{\mathsf{undef},\mathsf{void}\}\} \\
  P_{\mathsf{typeof}=\texttt{"object"}} &=
    \sem{\texttt{null}} \cup \{o \in \ObjVal \mid \neg\mathsf{callable}(o)\}\\
  P_{\mathsf{typeof}=\texttt{"function"}} &=
    \{o \in \ObjVal \mid \mathsf{callable}(o)\}
  \end{align*}
  and $P_{\mathsf{typeof}=k}$ the evident tag set for
  $k \in \{\texttt{"boolean"},\texttt{"number"},\texttt{"string"},
  \texttt{"bigint"},\texttt{"symbol"}\}$.
\item $\sigma \models_{\mathsf{true}} (\texttt{"}p\texttt{"}\ \texttt{in}\ x)$
  iff $\sigma(x) \in \ObjVal$ and $\mathsf{has}(\sigma(x), p)$;
  $\models_{\mathsf{false}}$ iff $\sigma(x) \in \ObjVal$ and
  $\neg\mathsf{has}(\sigma(x), p)$. (A non-object $\sigma(x)$ throws.)
\item $\sigma \models_{\mathsf{true}} (x\ \texttt{instanceof}\ C)$ iff
  $\sigma(x) \in P_{\mathsf{inst}}(C) \subseteq \ObjVal$
  (assumption \textbf{L-INST}).
\item $\sigma \models_{\mathsf{true}} \texttt{Array.isArray}(x)$ iff
  $\sigma(x) \in \ObjVal$ and $\mathsf{isArrayV}(\sigma(x))$;
  $\models_{\mathsf{false}}$ is the complement over all of $\Val$.
\item $\sigma \models_b\ (\texttt{!}e)$ iff $\sigma \models_{\neg b}\ e$.
\item $\sigma \models_{\mathsf{true}} (e_1\ \texttt{\&\&}\ e_2)$ iff
  $\sigma \models_{\mathsf{true}} e_1$ and
  $\sigma \models_{\mathsf{true}} e_2$;\;
  $\sigma \models_{\mathsf{false}} (e_1\ \texttt{\&\&}\ e_2)$ iff
  $\sigma \models_{\mathsf{false}} e_1$, or
  $\sigma \models_{\mathsf{true}} e_1$ and
  $\sigma \models_{\mathsf{false}} e_2$. Dually for \texttt{||}.
\item $\sigma \models_b\ x$ (truthiness of a variable) iff
  $\mathsf{truthy}(\sigma(x)) = b$; likewise
  $\sigma \models_b\ x.f$ iff $\sigma(x) \in \ObjVal$,
  $\mathsf{has}(\sigma(x), f)$ or the read yields \texttt{undefined}, and
  the read value's truthiness is $b$.
\end{itemize}
\end{definition}

\begin{definition}[Branch Soundness]
\label{def:branch-sound}
A rule deriving $\RefineB{\Env}{e}{\Env'}{b}$ is \emph{branch-sound} if
$\sigma \models_\rho \Env$ and $\sigma \models_b e$ imply
$\sigma \models_{\rho'} \Env'$, where $\rho \to \rho'$ is the substitution
evolution performed by the rule's refinement operators (an evolution of the
kind covered by \textbf{L-MONO}).
\end{definition}

\begin{lemma}[Predicate-Sound Update is Branch-Sound]
\label{lem:pred-to-branch}
If $R$ is predicate-sound for $P$, the update
$\Env' = \Env[x \mapsto R(\Env(x))]$ is branch-sound for any $e, b$ with
$\sigma \models_b e \implies \sigma(x) \in P$.
\end{lemma}
\begin{proof}
For unchanged variables $y$: $\sigma(y) \in \sem{\Env(y)}_\rho$, and
\textbf{L-MONO} carries this to $\rho'$ — this step is \emph{not} free, as
the semantics of $\Env(y)$ is state-indexed. For $x$:
$\sigma(x) \in \sem{\Env(x)}_\rho \cap P \subseteq \sem{R(\Env(x))}_{\rho'}$
by predicate-soundness of $R$.
\end{proof}

\subsection{Master list}
\label{sec:master-list}

Each row corresponds to \emph{one dispatch arm} of \texttt{evalCond}
(\texttt{TypeCheckRefine.sk}); the maintenance check of
Appendix~\ref{app:maintenance} keeps this correspondence mechanical. ``id''
means that control-flow branch performs no narrowing.

\begin{keydef}
\renewcommand{\arraystretch}{1.2}
\small
\begin{tabular}{@{}l@{\;\;}l@{\;\;}l@{\;\;}l@{}}
\toprule
\textbf{Rule} & \textbf{Condition form} & \textbf{Branches} & \textbf{Impl (dispatch arm / handler)} \\
\midrule
\textsc{R-Eq} & $x\ \texttt{===}\ e$, $x\ \texttt{===}\ y$ & both &
  \texttt{evalEquality}/\texttt{evalVarEquality} \\
\textsc{R-EqUndef} & $x\ \texttt{===}\ \texttt{undefined}$ & both &
  \texttt{evalVarEquality} (undefined case) \\
\textsc{R-EqField} & $x.f\ \texttt{===}\ e$ (and $x.f === y.g$) & both &
  \texttt{evalFieldEqualityType}, \texttt{narrowField} \\
\textsc{R-Typeof} & $\texttt{typeof}\ x\ \texttt{===}\ k$ & both &
  \texttt{evalTypeofEquality} \\
\textsc{R-IsArray} & \texttt{Array.isArray(}$x$\texttt{)} & both &
  \texttt{evalIsArray} \\
\textsc{R-Inst} & $x\ \texttt{instanceof}\ C$ & true only &
  \texttt{evalInstanceof} \\
\textsc{R-In} & $e_k\ \texttt{in}\ x$ & both &
  \texttt{evalCond} (\texttt{In} arm) \\
\textsc{R-Not} & $\texttt{!}e$ & both & \texttt{evalCond} (\texttt{Unop} arm)\\
\textsc{R-And} & $e_1\ \texttt{\&\&}\ e_2$ & both &
  \texttt{evalCond} (\texttt{BOLogicalAnd} arm) \\
\textsc{R-Or} & $e_1\ \texttt{||}\ e_2$ & both &
  \texttt{evalCond} (\texttt{BOLogicalOr} arm) \\
\textsc{R-Truthy} & $x$ in condition position & both &
  \texttt{evalTruthiness} \\
\textsc{R-TruthyField} & $x.f$ in condition position & both &
  \texttt{evalFieldTruthiness} \\
\textsc{R-Other} & anything else & both (id) &
  \texttt{evalCond} final wildcard \\
\bottomrule
\end{tabular}
\end{keydef}

\noindent
Entry points that reduce to this list: \texttt{if}/\texttt{while} conditions,
the ternary $c\ \texttt{?}\ a\ \texttt{:}\ b$ (which narrows inside $a$ and
$b$ exactly like \texttt{if}), and \texttt{\&\&}/\texttt{||} value positions
call \texttt{evalCond}; \texttt{switch} desugars each \texttt{case $v$:} to
$e\ \texttt{===}\ v$ (\texttt{TypeCheckStmt.sk});
\texttt{console.assert}$(e)$ applies the true branch of $e$ — matched by
call shape, first argument, without verifying \texttt{console} is the global
(\texttt{TypeCheckExpr.sk}). \texttt{for}-loop conditions do \emph{not}
narrow. Strict inequality dispatches to the same handlers with the target
negated.

\subsection{Rules and soundness}

Throughout, $b$ is the branch ($\mathsf{true}$ in the taken/then branch,
$\mathsf{false}$ in the untaken/else branch) and updates go through
Lemma~\ref{lem:pred-to-branch} unless noted.

\subsubsection{\textsc{R-Eq}: strict equality on a variable}
\label{sec:r-eq}

For $x\ \texttt{===}\ e$ with $e$ not the literal \texttt{undefined}
(that case is \textsc{R-EqUndef}), let $\tau_e$ be the checked type of $e$.
\begin{mathpar}
\inferrule*[right=R-EqT]
{ }
{\RefineT{\Env}{(x\,\texttt{===}\,e)}
  {\Env[x \mapsto \mathsf{inter}(\Env(x), \tau_e)]}}
\and
\inferrule*[right=R-EqF]
{\mathsf{singleConst}(\tau_e) = \mathsf{Some}(c)}
{\RefineF{\Env}{(x\,\texttt{===}\,e)}
  {\Env[x \mapsto \mathsf{rmc}(\Env(x), c)]}}
\end{mathpar}
where $\mathsf{rmc}$ is constant removal (Section~\ref{sec:removeconstant})
and $\mathsf{singleConst}$ succeeds only when $\tau_e$ is a single unit
or single-literal primitive, whose denotation is then \emph{contained in} a
single $\approx$-class — equal to it for literal primitives and the
\texttt{null} unit, a proper subset for the \texttt{undefined}/\texttt{void}
units, whose cross-tag partner
Definition~\ref{def:literal-denotation} deliberately excludes. If
$\mathsf{singleConst}$ fails, the false branch is the identity. When both
sides are variables (or a variable and a field path), the rule applies to
each side independently.
\impl{TypeCheckRefine.sk / evalVarEquality}

\begin{lemma}[\textsc{R-Eq} is Branch-Sound]
\label{lem:r-eq-sound}
Assume $\sem{\tau_e}$ is $\approx$-closed (assumption \textbf{L-SEQ}).
\end{lemma}
\begin{proof}
\emph{True branch.} $\sigma \models_{\mathsf{true}}$ gives
$\sigma(x) \approx v_e$, the value of $e$. By type correctness of the
checker for $e$ (assumption \textbf{L-TC}), $v_e \in \sem{\tau_e}$, and by
$\approx$-closure $\sigma(x) \in \sem{\tau_e}$. So
$\sigma(x) \in \sem{\Env(x)} \cap \sem{\tau_e} \subseteq
\sem{\mathsf{inter}(\Env(x), \tau_e)}$ by \textbf{L-INTER}.
As in Lemma~\ref{lem:pred-to-branch}, the unchanged variables transport to
the post-refinement state by \textbf{L-MONO} (this rule bypasses the
memberwise primitive, so the step must be taken explicitly).
\emph{False branch.} $\sigma(x) \not\approx v_e$, and every member of
$\sem{c}$ is ${\approx}\,v_e$ (Lemma~\ref{lem:per} for a literal constant,
since $v_e \in \sem{\tau_e} = \sem{c}$; unit reflexivity otherwise), so
$\sigma(x) \notin \sem{c}$; then
$\sigma(x) \in \sem{\Env(x)} \setminus \sem{c}
\subseteq \sem{\mathsf{rmc}(\Env(x), c)}$ by Theorem~\ref{thm:rmc}(R2).
\end{proof}

\paragraph{The $\approx$-closure hypothesis.}
It holds by construction for literal and full primitive types
(Definition~\ref{def:literal-denotation} takes the same-tag closure, using
transitivity via Lemma~\ref{lem:per}, and $\approx$ is payload equality on
non-numeric tags), for the \texttt{null} unit (payload equality at its own
tag — null comparisons flow through this generic rule), and trivially for
object types (reference equality). The corner where it can fail is the
cross-tag $\mathsf{undef}/\mathsf{void}$ arm of $\approx$, in either
direction: a $\tau_e$ denoting the \texttt{undefined} unit alone is not
closed against $\mathsf{void}$-tagged values, and dually for a
\texttt{void}-typed RHS meeting an undefined-tagged value. The syntactic
comparison
$x\ \texttt{===}\ \texttt{undefined}$ takes the dedicated rule
\textsc{R-EqUndef}, which treats the two tags jointly; a comparison against
an undefined- or void-\emph{typed} expression that is not the literal
(\texttt{x === y} with \texttt{y: undefined}) reaches the generic rule —
but the corner where either unit type could meet the other tag's value is
unreachable: the checker rejects \texttt{void}-typed operands of strict
equality outright (``An expression of type 'void' cannot be used here''),
before any refinement runs, discharging the last case of assumption
\textbf{L-SEQ} (empirically pinned).

\paragraph{State caveat.}
The implementation computes the intersection under a \emph{cloned}
substitution and installs the resulting type against the main one;
assumption \textbf{L-CLONE} records the stability requirement this
creates.

\subsubsection{\textsc{R-EqUndef}: comparison with \texttt{undefined}}

Mapper $g_b$: keep $s$ iff
$(\tagfn(s) \in \{\mathsf{undef}, \mathsf{void}\}) = b$; no default. This is
$\refineop^{g_b,\mathsf{None}}$ on both branches.
Soundness: $\sigma \models_b$ holds iff
$(\sigma(x) \approx (\mathsf{undef},\star)) = b$, and by
Definition~\ref{def:seq} that is exactly
$(\rtag(\sigma(x)) \in \{\mathsf{undef},\mathsf{void}\}) = b$ — the rule
treats the two tags jointly precisely because $\approx$ identifies them.
Apply Lemmas~\ref{lem:branchwise-sound} and~\ref{lem:pred-to-branch}.
(\texttt{null} comparisons are \emph{not} special-cased here: they flow
through \textsc{R-Eq} with $c = \texttt{null}$.)
\impl{TypeCheckRefine.sk / evalVarEquality (undefined case)}

\subsubsection{\textsc{R-EqField}: strict equality on a field path}
\label{sec:r-eqfield}

For $x.f\ \texttt{===}\ e$ with RHS type $\tau_e$ and
$c = \mathsf{singleConst}(\tau_e)$ (possibly undefined), the mapper acts only
on object branches. For each branch $\beta$ of each object member of
$\Env(x)$, let $\tau_\beta$ be $\beta$'s \emph{resolved read type} for the
\emph{declared named} field $f$: $\mathsf{get}(f)$ if $f$ is required,
$\mathsf{union}(\mathsf{get}(f), \texttt{undefined})$ if optional
(\texttt{maybeGetType}/\texttt{resolveFieldType} — the implementation
resolves before narrowing, which is load-bearing below). If $\beta$ has no
declared field $f$ — including when $f$ is reachable only through an index
signature — the object branch is kept unchanged (a sound over-approximation; no
diagnostic is emitted). Let $\omega = \mathsf{inter}(\tau_\beta, \tau_e)$:
\[
g_{\mathsf{true}}(\beta) =
\begin{cases}
\text{drop} & \omega = \bot\\
\beta[f.\mathsf{get} \mapsto \omega] & \text{else}
\end{cases}
\qquad
g_{\mathsf{false}}(\beta) =
\begin{cases}
\beta & \omega = \bot \text{ or } c\ \text{undefined}\\
\text{drop} & \mathsf{rmc}(\tau_\beta, c) = \bot\\
\beta[f.\mathsf{get} \mapsto \mathsf{rmc}(\tau_\beta, c)] & \text{else}
\end{cases}
\]
with the bookkeeping that narrowing out \texttt{undefined} also clears the
field's optionality.\footnote{The optionality-clearing test matches only the
\texttt{undefined} unit, not \texttt{void}; subtracting a
\texttt{void}-typed constant keeps optionality — conservative, hence sound.}
Non-object members are kept unchanged on both branches. The narrowed read
type is written back as the field's $\mathsf{get}$; when the field stays
optional, the next read re-adds \texttt{undefined}, which is conservative.
\impl{TypeCheckRefine.sk / evalFieldEqualityType, narrowField}

\begin{lemma}[\textsc{R-EqField} is Branch-Sound]
\label{lem:r-eqfield-sound}
Assume $\sem{\tau_e}$ is $\approx$-closed (\textbf{L-SEQ}), as in
\textsc{R-Eq}.
\end{lemma}
\begin{proof}
Let $v = \sigma(x) \in \sem{\Env(x)}$ and let the read value be
$w = \mathsf{get}(v, f)$ if present, else the value \texttt{undefined}.
For an object branch $\beta$ with declared field $f$, membership of $v$ in $\beta$
gives $w \in \sem{\tau_\beta}$ \emph{by construction of the resolved read
type}: if $f$ is present in $v$, Definition~\ref{def:object-semantics} gives
$w \in \sem{\mathsf{get}(f)} \subseteq \sem{\tau_\beta}$ (\textbf{L-UNION});
if $f$ is optional and absent,
$w = \texttt{undefined} \in \sem{\tau_\beta}$ because resolution unioned
\texttt{undefined} in. (This is why resolving before narrowing is
load-bearing: with the raw $\mathsf{get}$-type, the true branch of
\texttt{x.f === undefined} on an absent optional field would compute
$\omega = \bot$ and wrongly drop the object branch.)
\emph{True branch:} $w \approx v_e$ and $v_e \in \sem{\tau_e}$
(\textbf{L-TC}), so $w \in \sem{\tau_e}$ by closure, and
$w \in \sem{\omega}$ by \textbf{L-INTER}; in particular $\omega \neq \bot$,
so $v$'s object branch is not dropped. If $f$ is present in $v$, the strengthened
field is satisfied; if absent (possible only when
$v_e \approx \texttt{undefined}$), optionality is kept and the strengthened
field imposes no constraint.
\emph{False branch:} $w \not\approx v_e$. If $c$ is defined then
every member of $\sem{c}$ is ${\approx}\,v_e$ (as in \textsc{R-Eq}), so
$w \in \sem{\tau_\beta}\setminus\sem{c}
\subseteq \sem{\mathsf{rmc}(\tau_\beta, c)}$ by Theorem~\ref{thm:rmc}(R2);
since $w$ inhabits it, the removal is not $\bot$ and the object branch survives
with $v$ in it. When $c$ is the \texttt{undefined} unit, optionality is
cleared: $w \not\approx \texttt{undefined}$ entails $f$ is present in $v$,
so the now-required field is satisfied.
If $\omega = \bot$ then by \textbf{L-INTER} the field value could never
strictly equal $v_e$ (using closure of $\sem{\tau_e}$), and the object branch is
kept unchanged, which is sound. Branches without a declared $f$ and
non-object members are unchanged (identity,
Lemma~\ref{lem:identity-sound}).
\end{proof}

\subsubsection{\textsc{R-Typeof}}
\label{sec:r-typeof}

Eight recognized operands; any other string is a static error (TS2367), and
then no narrowing occurs. Writing $b$ for the control-flow branch:
\begin{itemize}
\item $k \in \{\texttt{"boolean"}, \texttt{"number"}, \texttt{"string"},
  \texttt{"bigint"}, \texttt{"symbol"}\}$: keep $s$ iff
  $(\tagfn(s) = \tagfn(k)) = b$; default on the true branch only:
  the full primitive of kind $k$.
\item $k = \texttt{"undefined"}$: keep $s$ iff
  $(\tagfn(s) \in \{\mathsf{undef},\mathsf{void}\}) = b$; no default.
\item $k = \texttt{"object"}$: on $b = \mathsf{true}$, keep \texttt{null};
  drop other primitives; per object branch, by call mode
  (Definition~\ref{def:sig-modes}): $\mathsf{None} \to$ keep,
  $(\mathsf{true},\sigma) \to$ keep with call mode $\mathsf{None}$,
  $(\mathsf{false},\sigma) \to$ drop.
  On $b = \mathsf{false}$: dual — \texttt{null} dropped, other primitives
  kept, per object branch $\mathsf{None} \to$ drop,
  $(\mathsf{true},\sigma) \to (\mathsf{false},\sigma)$,
  $(\mathsf{false},\sigma) \to$ keep. Default (true branch):
  $\{\texttt{null}\} \sqcup \mathsf{object}(\{\emptyset\})$.
\item $k = \texttt{"function"}$: on $b = \mathsf{true}$: drop all
  non-objects; per object branch $\mathsf{None} \to$ drop,
  $(\mathsf{true},\sigma) \to (\mathsf{false},\sigma)$,
  $(\mathsf{false},\sigma) \to$ keep. On $b = \mathsf{false}$: keep
  non-objects; per object branch $\mathsf{None} \to$ keep,
  $(\mathsf{true},\sigma) \to$ call mode $\mathsf{None}$,
  $(\mathsf{false},\sigma) \to$ drop. No default.
\end{itemize}
\impl{TypeCheckRefine.sk / evalTypeofEquality}

\begin{lemma}[\textsc{R-Typeof} is Branch-Sound]
\label{lem:r-typeof-sound}
\end{lemma}
\begin{proof}
Tag cases are immediate from A0 and Definition~\ref{def:eval}; defaults are
adequate since $P_{\mathsf{typeof}=k} \subseteq \sem{\mathrm{full}\ k}$. For
\texttt{"object"}, the default's object branch (\texttt{TObject::createEmpty})
has no fields or indexes and $\mathsf{call} = \mathsf{new} = \mathsf{None}$,
which by Definition~\ref{def:object-semantics} constrains members to be
non-callable and non-constructible; every
$o \in P_{\mathsf{typeof}=\texttt{"object"}} \setminus \sem{\texttt{null}}$
is non-callable by definition and therefore non-constructible by the
$\mathsf{constructible} \implies \mathsf{callable}$ fact of
Definition~\ref{def:object-interface}, so adequacy holds. For the callable cases, let
$v \in \ObjVal \cap \sem{\beta}$. If $\mathsf{callable}(v)$
($k=\texttt{"function"}$, $b=\mathsf{true}$): $\beta$'s call mode cannot be
$\mathsf{None}$ (Definition~\ref{def:object-semantics}); a
$(\mathsf{true},\sigma)$ mode is strengthened to
$(\mathsf{false},\sigma)$, which $v$ satisfies since it is callable.
If $\neg\mathsf{callable}(v)$: the mode cannot be $(\mathsf{false},\sigma)$;
a $(\mathsf{true},\sigma)$ mode weakens/strengthens to $\mathsf{None}$, which
$v$ satisfies. The \texttt{"object"} cases combine this with the
\texttt{null} tag case. The mode rewrites preserve all other constraints of
$\beta$, and are tag-preserving, so Lemma~\ref{lem:branchwise-sound} applies.
\end{proof}

\subsubsection{\textsc{R-IsArray}}
\label{sec:r-isarray}

True branch: drop non-object members; each object branch is intersected with
the single object branch of $\mathsf{object}(\mathsf{ReadonlyArray}\langle\top\rangle)$ and
its brand set to $\mathsf{Some}(\mathsf{true})$; an object branch whose intersection
is $\bot$, or that was already branded $\mathsf{Some}(\mathsf{false})$, is
dropped. (The \texttt{ReadonlyArray} stub itself carries the brand as a
declared field \texttt{readonly \_\_SKJS\_isArray:\ true}, so the
brand-conflict drop already occurs inside the intersection.) False branch:
object branches have the brand set to $\mathsf{Some}(\mathsf{false})$; a
branch already branded $\mathsf{Some}(\mathsf{true})$ is dropped;
non-objects are kept.
\impl{TypeCheckRefine.sk / evalIsArray; TTypeInter.sk /
TFields::interIsArray}

\begin{lemma}[\textsc{R-IsArray} is Branch-Sound]
\label{lem:r-isarray-sound}
\end{lemma}
\begin{proof}
True branch: $\sigma \models_{\mathsf{true}}$ gives
$v = \sigma(x) \in \ObjVal$ with $\mathsf{isArrayV}(v)$. By assumption
\textbf{L-ISARR}, $v \in
\sem{\mathsf{object}(\mathsf{ReadonlyArray}\langle\top\rangle)}$; by
\textbf{L-INTER} (object level) $v$ inhabits the intersected branch, and
setting the brand to $\mathsf{Some}(\mathsf{true})$ is satisfied by
$\mathsf{isArrayV}(v)$. False branch: if $v \in \ObjVal$ then
$\neg\mathsf{isArrayV}(v)$ satisfies brand $\mathsf{Some}(\mathsf{false})$
and contradicts membership in a $\mathsf{Some}(\mathsf{true})$-branded
branch, so dropping those loses nothing; non-objects pass through.
\end{proof}

\paragraph{Mutability note.}
Intersecting with \emph{Readonly}Array rather than \texttt{Array} is a
deliberate precision/soundness choice documented in the implementation: a
readonly-array-typed value may be a non-array object carrying the right
properties in general, but under $\mathsf{isArrayV}$ the brand is known; what
is \emph{not} known is mutability, so the mutable interface is added only
where the original branch already had it.

\subsubsection{\textsc{R-Inst}: \texttt{instanceof}}
\label{sec:r-inst}

Fires only on the true branch and only when the RHS names a class or
interface with a known definition (a Function-typed value as RHS narrows
nothing). Let $T_C = \mathsf{object}$-typed instance of $C$ with
implementation-chosen type arguments ($\top$ for \texttt{Array}; fresh
inference arguments otherwise). Mapper: drop non-objects; intersect each
object member with $T_C$ (dropping empty intersections). Default:
$T_C$ itself.
\impl{TypeCheckRefine.sk / evalInstanceof, evalCond (InstanceOf arm)}

\begin{lemma}[\textsc{R-Inst} is Branch-Sound]
Given assumption \textbf{L-INST}
($P_{\mathsf{inst}}(C) \subseteq \sem{T_C}$, with $T_C$ as instantiated at
the narrowing site), the true-branch rule is branch-sound; the false branch
is the identity.
\end{lemma}
\begin{proof}
Let $v = \sigma(x)$ with $\sigma \models_{\mathsf{true}}$, so
$v \in P_{\mathsf{inst}}(C) \subseteq \ObjVal$. Dropping non-object members
is sound by A0 (values of a non-object member carry that member's tag, so
$v$ is not in it). For $v$'s object member,
$v \in \sem{T_C}$ by \textbf{L-INST}, so
$v$ inhabits the intersection by \textbf{L-INTER} (object level), which is
in particular not $\bot$. The default $D = \{T_C\}$ is a singleton
$\STypes$ ($T_C$ is a \texttt{TObject} built by \texttt{newObject}) and is
adequate by \textbf{L-INST} directly. Both \textbf{L-INST} and the mapper
are evaluated at the post-instantiation state $\rho'$ ($T_C$'s fresh
inference arguments live in $\rho'$); later growth of those arguments is
covered by \textbf{L-MONO}/Obligation~\ref{obl:monotone}. Apply
Lemmas~\ref{lem:branchwise-sound} and~\ref{lem:pred-to-branch}.
\end{proof}

\subsubsection{\textsc{R-In}}
\label{sec:r-in}

The key expression's \emph{type} determines the refinement: if it denotes a
single string \emph{or numeric} literal $p$ or a single (unique) symbol
$q$, narrowing proceeds with key $\pi \in \{p, q\}$ — a numeric literal is
converted to its string spelling and the key kind is hardcoded to string,
so the true branch's index lookup below consults the \emph{string} index
even for a numeric key — sound under the covering semantics of
Definition~\ref{def:object-semantics}, since a string index covers
numeric-kind keys too; the imprecision is only that a number index is not
consulted. If the key type is a multi-literal or non-literal
string/number/symbol type, object members are kept unchanged — but
primitive members are still dropped on the true branch; if it is not
a string/number/symbol type at all, the checker raises a \emph{fatal}
TS2360 and no narrowing (or further checking of the definition) occurs.
In every non-fatal case the true branch drops primitive members (evaluation
of \texttt{in} on a primitive throws; Definition~\ref{def:eval}). With a
single key $\pi$, per object branch $\beta$:
\begin{itemize}
\item \textbf{True branch:} if $\beta$ has field $\pi$: clear its
  optionality. Else if $\beta$ has a non-synthetic index \emph{covering} $\pi$
  with field $f_{\mathrm{idx}}$: add the required field
  $\pi \mapsto f_{\mathrm{idx}}$. Else: add a required \emph{phantom} field
  $\pi$ ($\mathsf{get} = \top$, $\mathsf{set} = \bot$,
  $\mathsf{access} = \mathsf{Phantom}$).
\item \textbf{False branch:} if $\beta$ has required $\pi$: drop the object branch.
  If optional $\pi$: remove the field from $\beta$. Else: keep.
\end{itemize}
There is no default: a top-clause type stays top on both control-flow
branches.
\impl{TypeCheckRefine.sk / evalCond (In arm), inLhsFieldsFromType}

\begin{lemma}[\textsc{R-In} is Branch-Sound]
\label{lem:r-in-sound}
\end{lemma}
\begin{proof}
True branch: $v = \sigma(x) \in \ObjVal$ with $\mathsf{has}(v, \pi)$.
Primitives cannot contain $v$ (A0), so dropping them is sound. If $\beta$
had optional $\pi$, presence upgrades $v$'s satisfaction to the required
form; the value constraint is unchanged. If the field is added from an index
covering $\pi$: membership of $v$ in $\beta$ already constrains
$\mathsf{get}(v,\pi) \in \sem{f_{\mathrm{idx}}.\mathsf{get}}$
(Definition~\ref{def:object-semantics}, index clause — $\pi$ is covered
and was not named), so $v$ satisfies the named copy. A phantom field constrains presence only, which holds.
False branch: $v \in \ObjVal$ with $\neg\mathsf{has}(v,\pi)$ contradicts a
required $\pi$ — including a required \emph{phantom}, which still demands
presence — so $v$ was never in dropped object branches. Removing an optional field
$\pi$ is \emph{not} monotone in general (the removal re-exposes $\pi$ to
the object branch's index constraint, which the named field had exempted), but it
is sound here: $v$ lacks $\pi$, so the re-exposed index constraint on
$\pi$ is vacuous for $v$, and constraints on other keys are unchanged.
Primitives kept is identity.
\end{proof}

\subsubsection{\textsc{R-Not}, \textsc{R-And}, \textsc{R-Or}}
\label{sec:connectives}

\begin{mathpar}
\inferrule*[right=R-Not]
{\RefineB{\Env}{e}{\Env'}{\neg b}}
{\RefineB{\Env}{\texttt{!}e}{\Env'}{b}}
\and
\inferrule*[right=R-AndT]
{\RefineT{\Env}{e_1}{\Env_1} \\ \RefineT{\Env_1}{e_2}{\Env_2}}
{\RefineT{\Env}{(e_1\,\texttt{\&\&}\,e_2)}{\Env_2}}
\and
\inferrule*[right=R-AndF]
{\RefineF{\Env}{e_1}{\Env_1} \\
 \RefineT{\Env}{e_1}{\Env'} \\ \RefineF{\Env'}{e_2}{\Env_2}}
{\RefineF{\Env}{(e_1\,\texttt{\&\&}\,e_2)}{\Env_1 \sqcup \Env_2}}
\end{mathpar}
and dually \textsc{R-OrF} (sequential) and \textsc{R-OrT} (merge of
``$e_1$ true'' and ``$e_1$ false then $e_2$ true''). Here
$(\Env_1 \sqcup \Env_2)(x) = \mathsf{union}(\Env_1(x), \Env_2(x))$ is the
pointwise scope merge. (The implementation merges only the bindings the
sub-scopes updated: a variable narrowed in just one sub-scope is unioned
with the \emph{parent} scope's type — which equals the pointwise union,
since the un-narrowed side's value is the parent's.)
\impl{TypeCheckRefine.sk / evalCond (Unop / BOLogicalAnd / BOLogicalOr
arms);
TypeCheckScope.sk / Scope::merge via TypeCheckEnv.sk / mergeScope}

\begin{lemma}[Merge Soundness]
\label{lem:merge-sound}
If $\sigma \models \Env_1$ or $\sigma \models \Env_2$, then
$\sigma \models \Env_1 \sqcup \Env_2$.
\end{lemma}
\begin{proof}
Pointwise from \textbf{L-UNION}. The domains agree
($\mathrm{dom}(\Env_1) = \mathrm{dom}(\Env_2)$) because every rule is an
update $\Env[x \mapsto \cdot]$ of an existing local — refinement never adds
or removes bindings. The merged environment, like all post-refinement
environments, is interpreted at the final state; the premises' satisfaction
transports to it by \textbf{L-MONO}.
\end{proof}

\begin{lemma}[Connectives are Branch-Sound]
\label{lem:connectives-sound}
If all premise derivations are branch-sound, so are
\textsc{R-Not}, \textsc{R-And}, \textsc{R-Or} — for \emph{both} branches.
\end{lemma}
\begin{proof}
\textsc{R-Not} is immediate from Definition~\ref{def:eval}.
\textsc{R-AndT} composes directly: both conjuncts hold and the two
refinements chain. For \textsc{R-AndF}: by Definition~\ref{def:eval},
$\sigma \models_{\mathsf{false}} e_1$ — giving $\sigma \models \Env_1$ — or
$\sigma \models_{\mathsf{true}} e_1$ and
$\sigma \models_{\mathsf{false}} e_2$ — giving $\sigma \models \Env'$ then
$\sigma \models \Env_2$. Either way Lemma~\ref{lem:merge-sound} concludes.
The \texttt{||} cases are dual. Note the exactness of the short-circuit
semantics in Definition~\ref{def:eval} is what licenses refining
\emph{both} control-flow branches of both connectives.
\end{proof}

\subsubsection{\textsc{R-Truthy}, \textsc{R-TruthyField}}
\label{sec:truthiness}

\begin{definition}[Truthiness Residue]
\label{def:residue}
For $b \in \{\mathsf{true},\mathsf{false}\}$, the partial mapper
$\mathsf{tres}_b : \SType \rightharpoonup \SType$ is:
\begin{itemize}
\item $\mathsf{unit}(\texttt{void})$: kept unchanged on both branches
  (\texttt{void} is treated as non-observable and never narrowed);
\item $\mathsf{unit}(\texttt{undefined}), \mathsf{unit}(\texttt{null})$:
  dropped if $b$, kept if $\neg b$;
\item objects and symbols: kept if $b$, dropped if $\neg b$;
\item full \texttt{boolean}: the literal $b$;
\item full \texttt{string}/\texttt{number}/\texttt{bigint}: kept unchanged
  on \emph{both} branches;
\item literal primitives $\mathsf{prim}(k, L)$: keep the literals
  $\ell \in L$ with $\mathsf{truthy}(\delta_k(\ell)) = b$; undefined if none
  remain.
\end{itemize}
The \emph{value} residue $\mathsf{vres}_b$ differs in exactly two cases:
full \texttt{string} and full \texttt{bigint} on $b = \mathsf{false}$ narrow
to the literals \texttt{""} and \texttt{0n} (their unique falsy values);
full \texttt{number} must stay wide because $\mathsf{NaN}$ has no literal
type.
\impl{TType.sk / SType::truthinessResidue, truthinessValueResidue}
\end{definition}

\begin{lemma}[Residues are Sound]
\label{lem:residue-sound}
$\mathsf{tres}_b$ is sound (Definition~\ref{def:pred-sound}) for
$P_b = \{v \mid \mathsf{truthy}(v) = b\} \cup \sem{\texttt{void}}$, and
$\mathsf{vres}_b$ is sound for $P_b$ as well.
\end{lemma}
\begin{proof}
Case analysis over Definition~\ref{def:truthiness}. The two deliberate
widenings (full string/number/bigint kept on both branches) only enlarge the
result. The literal-filtering case is sound \emph{only if} the
implementation's lexeme-level falsiness test agrees with
Definition~\ref{def:truthiness} on every literal in $L$ — this is the
falsiness-exactness part (iii) of assumption \textbf{L-LIT}, which the
string kind violates (see the assumptions table for the full
inventory of violated assumptions): a line-continuation string lexeme
denotes \texttt{""} yet is classified truthy, and would be dropped from
the falsy branch. For $\mathsf{vres}_{\mathsf{false}}$: the only falsy string value is
$\texttt{""}$ and the only falsy bigint value is $0$, so the narrowed
literal contains every falsy member; number stays wide precisely because its
falsy set $\{0, -0, \mathsf{NaN}\}$ is not expressible as literals.
The $\sem{\texttt{void}}$ disjunct in $P_b$ accounts for keeping
\texttt{void} on both branches (for $b = \mathsf{false}$ it is redundant,
since void values are falsy).
\end{proof}

\textsc{R-Truthy} is $\refineop^{\mathsf{tres}_b, \mathsf{None}}$ on
$\Env(x)$; branch soundness follows from
Lemmas~\ref{lem:branchwise-sound}, \ref{lem:residue-sound}
and~\ref{lem:pred-to-branch} (note $\sigma \models_b x$ implies
$\sigma(x) \in P_b$, and \texttt{void}-typed reads evaluate to
\texttt{undefined}, which the kept \texttt{void} member covers).
\textsc{R-TruthyField} applies the \emph{type-level} residue
($\refineop^{\mathsf{tres}_b,\mathsf{None}}$ on the field's type) through the
same per-branch field update as \textsc{R-EqField}, dropping object branches whose
residue is $\bot$ and clearing optionality on the true branch (the read
value being truthy entails the field is present with a non-\texttt{undefined}
value); its proof combines Lemma~\ref{lem:r-eqfield-sound}'s structure with
Lemma~\ref{lem:residue-sound}.
\impl{TypeCheckRefine.sk / evalTruthiness, evalFieldTruthiness;
TType.sk / TType::truthinessPart}

\paragraph{Value positions of \texttt{\&\&}/\texttt{||}.}
The \emph{type of the expression} $e_1\,\texttt{\&\&}\,e_2$ uses
$\mathsf{vres}_{\mathsf{false}}$ on $e_1$'s type for the untaken-operand
contribution (and dually for \texttt{||}); this is a typing rule rather than
a refinement and inherits soundness from Lemma~\ref{lem:residue-sound}.
\impl{TType.sk / TType::truthinessValuePart}

%==========================================================================
\section{Constant Removal}
\label{sec:removeconstant}

SKJS has no general type-difference operation: inference growth is
structural (Section~\ref{sec:inference}) and consumes no differences, and
the only subtraction clients — destructuring defaults
($\tau \setminus \texttt{undefined}$) and \texttt{!==} narrowing —
subtract a \emph{single constant}. (The historical \texttt{TType::diff}
and its per-level strata were deleted when growth became structural.)

\begin{definition}[Constant Removal]
\label{def:rmc}
A \emph{constant} $c$ is a unit or a single-literal primitive; its
denotation $\sem{c}$ is \emph{contained in} a single $\approx$-class —
equal to it for literal primitives (for a numeric constant this may be the
two-element class $\{0, -0\}$) and the \texttt{null} unit; a proper
subset for the \texttt{undefined}/\texttt{void} units, whose cross-tag
partner Definition~\ref{def:literal-denotation} deliberately excludes.
$\mathsf{rmc}(t, c)$ expands $t$ and maps each clause's $\STypes$ elementwise:
\begin{itemize}
\item $\mathsf{unit}(u)$ with $\tagfn(u) = \tagfn(c)$: dropped;
\item $\mathsf{prim}(k, L)$ with $\tagfn(k) = \tagfn(c)$ and $L \neq \top$:
  becomes $\mathsf{prim}(k, L \setminus \{\ell_c\})$, dropped if empty;
\item full \texttt{boolean} minus a boolean literal: the complementary
  literal (exact);
\item full \texttt{number}/\texttt{string}/\texttt{bigint}/\texttt{symbol}:
  kept unchanged (over-approximation);
\item all other members (different tag, objects): kept unchanged.
\end{itemize}
Top clauses are unchanged. The instance layer is unchanged.
\impl{TType.sk / TType::removeConstant; PrimLiterals.sk / diff}
\end{definition}

\begin{theorem}[Soundness of Constant Removal]
\label{thm:rmc}
For all $t, c$ (and states, suppressed):
\begin{itemize}
\item[(R1)] $\sem{\mathsf{rmc}(t,c)} \subseteq \sem{\mathsf{expand}(t)}$;
\item[(R2)] $\sem{t} \setminus \sem{c} \subseteq \sem{\mathsf{rmc}(t,c)}$.
\end{itemize}
Moreover removal is \emph{exact}
($\sem{\mathsf{rmc}(t,c)} = \sem{\mathsf{expand}(t)} \setminus \sem{c}$
on the clause part) whenever the expansion has no top clause and every
same-tag member is a unit, a literal primitive, or full \texttt{boolean} —
assuming additionally the injectivity part (ii) of \textbf{L-LIT}
(distinct normalized lexemes denote $\approx$-inequivalent values).
\end{theorem}
\begin{proof}
(R1): every output member's semantics is a subset of the corresponding input
member's ($L \setminus \{\ell_c\} \subseteq L$; the boolean complement is a
subset of full boolean; kept members are equal).
(R2): let $v \in \sem{t} \setminus \sem{c}$; by \textbf{L-EXP}
$v \in \sem{\mathsf{expand}(t)}$. If $v$ is in the instance layer or a top
clause, it survives unchanged. Otherwise $v \in \sem{s}$ for a member $s$.
If $\tagfn(s) \neq \tagfn(c)$, $s$ is kept. If $s$ is a unit of $c$'s tag,
then $s = c$ (the tag determines the unit) so $\sem{s} = \sem{c}$,
contradicting $v \notin \sem{c}$ — vacuous. (A \texttt{void}-tagged $v$
with $c$ the \texttt{undefined} unit falls under the \emph{different-tag}
case: $\mathsf{unit}(\texttt{void})$ is kept, and
$(\mathsf{void},\star) \notin \sem{c}$ under the singleton unit
denotations, so R2 is unaffected.) If $s$ is a literal primitive,
$v \approx (\tagfn(k), \delta_k(\ell))$ for some $\ell \in L$; if
$\ell = \ell_c$ then $v \in \sem{c}$, contradiction — so $\ell \neq \ell_c$
and $\ell$ survives the lexeme diff, keeping $v$ in the closure semantics.
Note this step needs only that $\delta$ is a \emph{function}, not
injectivity: R2 is unconditional. If $s$ is full boolean and $c$ a boolean
literal, $v$ is the complementary literal. Full other primitives are kept
whole. Exactness: a top clause is unchanged while
$\sem{\top} \setminus \sem{c} \subsetneq \sem{\top}$, hence its exclusion;
in the remaining listed cases each step removes exactly $\sem{c}$, using
injectivity to ensure no \emph{other} surviving lexeme denotes a member of
$\sem{c}$. (R1 has no client among the soundness results in this document;
its role is precision.)
\end{proof}

\paragraph{Subtracting from top.}
A top clause is left unchanged: \texttt{unknown} minus a constant is still
\texttt{unknown}. Narrowing \texttt{unknown} by
\texttt{!== null \&\& !== undefined} therefore does not narrow at all —
an incompleteness localized to the top-clause case of
Definition~\ref{def:rmc} and to the residue rules (a top clause has no
members to filter). Expressing a ``non-nullish top''
would still require a representation extension.

%==========================================================================
\section{Inference: Substitution Growth and Validation}
\label{sec:inference}

This section models the substitution state machine exactly and states the
soundness architecture as explicit obligations.

\subsection{The substitution state machine}

\begin{definition}[Substitution States]
\label{def:subst}
A substitution state $\rho$ maps type-variable identifiers to
$\SVOpen(t)$ (a growable lower bound) or $\SVClosed(t)$ (a sealed value),
and carries a boolean $\mathsf{freeze}$ (the candidate-freezing flag).
Exactly the \emph{ordinary} identifiers (\texttt{TVID}) are
\emph{candidates} (\path{TVarID::isCandidate}); field-inference
(\texttt{TVIDForField}) and literal-commitment
(\path{TVIDForLiteralCommitment}) identifiers are \emph{not} — the code
classifies both as literal-form commitments, and they \emph{are} sealed or
widened by a covariant grow. Lemma~\ref{lem:never-widen} therefore protects
ordinary inference candidates only.
\impl{TTypeSubst.sk / SubstValue, Subst; TType.sk / TVarID::isCandidate}
\end{definition}

\begin{definition}[Transitions]
\label{def:transitions}
$\mathsf{grow}(\alpha, t, \mathit{var})$ and
$\mathsf{get}(\alpha, \mathit{var})$ act as follows
(writing $\vee$ for \texttt{TType::union}):
\begin{center}
\renewcommand{\arraystretch}{1.2}
\begin{tabular}{@{}llll@{}}
\toprule
op & state of $\alpha$ & condition & effect / result \\
\midrule
$\mathsf{grow}(t,\mathsf{contra})$ & unbound & &
  $\alpha \mapsto \SVOpen(t)$ \\
 & $\SVOpen(t_2)$ & $\neg(\mathsf{cand}(\alpha) \wedge \mathsf{freeze})$ &
  $\alpha \mapsto \SVOpen(t_2 \vee t)$ \\
 & $\SVOpen(t_2)$ & $\mathsf{cand}(\alpha) \wedge \mathsf{freeze}$ & no-op \\
 & $\SVClosed(\_)$ & & no-op \\
\midrule
$\mathsf{grow}(t,\mathsf{cov})$ & unbound & &
  $\alpha \mapsto \SVClosed(t)$ \\
 & $\SVOpen(t_2)$ & $\mathsf{cand}(\alpha)$ & no-op \emph{(unconditional)} \\
 & $\SVOpen(t_2)$ & $\neg\mathsf{cand}(\alpha)$ &
  $\alpha \mapsto \SVClosed(t_2 \vee t)$ \\
 & $\SVClosed(\_)$ & & no-op \\
\midrule
$\mathsf{get}(\mathsf{cov})$ & unbound & & seal $\alpha \mapsto
  \SVClosed(\top)$; return $\top$ \\
 & $\SVOpen(t)$ & $\mathsf{cand}(\alpha) \wedge \mathsf{freeze}$ &
  return $t$ (no seal) \\
 & $\SVOpen(t)$ & otherwise & seal $\alpha \mapsto \SVClosed(t)$; return $t$\\
 & $\SVClosed(t)$ & & return $t$ \\
\midrule
$\mathsf{get}(\mathsf{contra})$ & unbound & & return $\bot$ (no write) \\
 & $\SVOpen(t)$ or $\SVClosed(t)$ & & return $t$ \\
\bottomrule
\end{tabular}
\end{center}
Both $\mathsf{get}$s then intersect the result with $\alpha$'s constraint
types. \impl{TTypeSubst.sk / Subst::grow, Subst::get}
\end{definition}

\begin{lemma}[Determined Candidates Are Never Widened]
\label{lem:never-widen}
For a candidate $\alpha$ bound to $\SVOpen(v)$ or $\SVClosed(v)$:
(i) a covariant $\mathsf{grow}$ never changes $v$;
(ii) under $\mathsf{freeze}$, no $\mathsf{grow}$ or $\mathsf{get}$ changes
$v$ or its open/closed status;
(iii) once $\SVClosed$, the binding never changes again.
\end{lemma}
\begin{proof}
By inspection of Definition~\ref{def:transitions}: (i) the
$\SVOpen$/candidate/covariant row is an unconditional no-op and the
$\SVClosed$ rows are no-ops; (ii) the two $\mathsf{freeze}$-guarded rows are
the only remaining writes touching a bound candidate, and both are
suppressed; (iii) no transition rewrites an $\SVClosed$ binding.
\end{proof}

This is the write-site invariant behind the
``covariant checks never widen a determined inference candidate'' work
(\#713, \#740): a covariant grow pushes an \emph{expected} type — an upper
bound — and the value of a candidate must come from its defining expression,
with the comparison against the expected type happening in validation, not by
overwriting the candidate.

\paragraph{Write-site audit.}
The lemma's premise — that Definition~\ref{def:transitions} captures
\emph{every} write to the binding map — has been checked against the source:
the only assignments to \texttt{Subst.values} are in \texttt{grow},
\texttt{get}, and the bulk \texttt{restore}. External callers of
\texttt{grow} (array-literal commitment in \texttt{TArrayLitTyping.sk},
variadic and utility sites in \texttt{TVariadicTyping.sk} /
\texttt{TypeCheckUtils.sk}) go through the same transitions;
\texttt{TVar::setConstraint} (used by intersection) is a \emph{functional}
update returning a new \texttt{TVar} with an extended constraint set —
constraints are occurrence-level, not shared state (registry row
SETCONSTRAINT-FUN) — and is
accounted for by the constraint intersection in $\mathsf{get}$.

\subsection{Growth as pre-seeding; validation decides}

\begin{definition}[Subtyping Check Architecture]
\label{def:issub-arch}
\texttt{isSub}$(t_1, t_2)$ proceeds as:
(0) acceptance shortcuts that return before any growth --- $t_1 = t_2$, and
the reflexive \texttt{TSInst} arm (equal name, arguments equal after
\texttt{peekRaw}), both guidance-free; the \texttt{TSInst} arm's
\texttt{peekRaw} reads an open candidate \emph{without sealing it},
the unsealed-read instance of
Obligation~\ref{obl:monotone} (deep dive D);
(1) $\mathsf{growTVars}(t_1, t_2, \mathsf{contra})$ — grow variables of the
to-be-grown side $t_2$ (the implementation's ``open'' side) from bound
material in $t_1$;
(2) $\mathsf{growTVars}(t_2, t_1, \mathsf{cov})$ — seal variables of $t_1$;
(3) expand both sides under the (now grown) state and validate with
\texttt{Clauses::isSub}; on failure, \emph{restore} the substitution to its
state before (1).
$\mathsf{growTVars}$ itself: collects growable variables of the
to-be-grown side
(bare, inside intersections if unbound, and behind one-level alias
expansion); checks each variable's declared constraints against the bound
side \emph{under} $\mathsf{freeze}$ (a query, per
Lemma~\ref{lem:never-widen}(ii)), rolling back on failure; grows only the
disjuncts of the bound side not already covered by a non-variable branch of
that side (each coverage probe is save/restored); finally descends
structurally into matching object fields and call/new \emph{return} types.
\impl{TTypeSubtyping.sk / isSub\_, growTVars, collectTVars}
\end{definition}

\begin{theorem}[Guidance vs.\ Validation]
\label{thm:guidance}
Assume the validation obligation \textbf{L-SUB}: if
\texttt{Clauses::isSub} accepts under state $\rho$, then
$\sem{t_1}_\rho \subseteq \sem{t_2}_\rho$ for the expanded views. Then no
imprecision in $\mathsf{growTVars}$ — growing too little, too much, or in
the wrong order — can cause \texttt{isSub} to accept an inclusion that is
invalid \emph{under the state $\rho'$ that growth produced}. Growth affects
only \emph{completeness} (which valid programs are accepted) and the
\emph{solutions} chosen for inference variables. (That already-validated
judgements are not retroactively weakened by later evolution is a separate
matter: Lemma~\ref{lem:never-widen} for candidates,
Obligations~\ref{obl:monotone} and~\ref{obl:restore} in general; the
failure-path parenthetical in the proof leans on
Obligation~\ref{obl:restore}.)
\end{theorem}
\begin{proof}
Acceptance occurs only in step (3), whose verdict is by \textbf{L-SUB} a
semantic inclusion \emph{under the state that growth produced}. Growth
selects the state; it never bypasses validation. (Failure paths restore the
prior state, so failed exploration does not corrupt later checks — see
Obligation~\ref{obl:restore}.)
\end{proof}

The force of Theorem~\ref{thm:guidance} is architectural rather than
representational. What it makes explicit is that the meaning of
\emph{other, already-checked} judgements must not be retroactively weakened
by later growth. That is exactly what
Lemma~\ref{lem:never-widen} secures for candidates, and what the remaining
obligations must secure in general:

\begin{obligation}[Validation Soundness, \textbf{L-SUB}]
\label{obl:sub}
\texttt{Clauses::isSub} (with its $\STypes$, object-branch, field,
index, and function strata, including getter/setter variance and the
\texttt{seenSub}/\texttt{tenv.seenSub} cycle caches keyed by substitution
version) is sound in the sense of Definition~\ref{def:issub-arch}. This is
the largest unproved component of the document.
\impl{TTypeSubtyping.sk / Clauses::isSub, STypes::isSub, TObject::isSub,
TFields::isSub, TFun::isSub}
\end{obligation}

\begin{obligation}[Monotone Reads Under Growth (registry row
OBL-MONOTONE)]
\label{obl:monotone}
If a judgement was validated under $\rho$ and inference later evolves
$\rho \to \rho'$ by the transitions of Definition~\ref{def:transitions},
the validated inclusion still holds under $\rho'$. For $\SVClosed$ bindings
this is Lemma~\ref{lem:never-widen}(iii); for $\SVOpen$ bindings of
\emph{non-candidates}, unions grow the denotation and the obligation is that
every consumer read them contravariantly (as lower bounds) until sealed.
\emph{This obligation is \textsc{violated} in
two witnessed instances} — \texttt{subst.seenSub} memos recorded
on unsealed candidate bounds inside \texttt{freezeCandidates} windows
(order-dependent acceptance), and the reflexive \texttt{TSInst} arm of
Definition~\ref{def:issub-arch} step (0), whose \texttt{peekRaw} is an
exact-value read of an unsealed open candidate (not memo-mediated). The
general statement is proved by deep dive
E modulo exactly these residues; witnesses and fix directions:
\path{deep-dives/issues/freeze-window-stale-memo-caches.md},
\path{deep-dives/issues/reflexive-tsinst-fastpath-unsealed-read.md}.
\end{obligation}

\begin{obligation}[No Commitment on Failure (registry row
OBL-RESTORE)]
\label{obl:restore}
Every failing or exploratory path restores the substitution:
\texttt{isSub} restores on validation failure; coverage probes and
constraint-check failures inside $\mathsf{growTVars}$ restore; refinement's
exploratory intersections run on clones (\textbf{L-CLONE}). The stale-cache
regression class (\#446) is guarded by keying cycle caches with the
substitution version.
\end{obligation}

%==========================================================================
\section{Worked Examples}
\label{sec:examples}

\subsection{Literal narrowing}
A worked case of \textsc{R-Eq}'s false branch:
\begin{verbatim}
function f(x: 1 | 2): 2 {
  if (x === 1) { return 2; }
  return x;  // narrowed to 2
}
\end{verbatim}
is accepted: \texttt{1 | 2} is
$\mathsf{prim}(\texttt{num}, \{1,2\})$; the false branch of \textsc{R-Eq}
applies $\mathsf{rmc}$ with $c = 1$, and Theorem~\ref{thm:rmc} gives exactly
$\mathsf{prim}(\texttt{num}, \{2\})$.

\subsection{Persisting incompleteness: subtracting from top}
\texttt{unknown} is the top clause; \textsc{R-Eq}'s false branch and the
truthiness residues leave top clauses unchanged
(Section~\ref{sec:removeconstant}), so
\texttt{x !== null \&\& x !== undefined} does not narrow
\texttt{x: unknown}. This remains a representation limitation: there is no
form for ``all values except these tags''.

\subsection{Persisting incompleteness: no false-branch \texttt{instanceof}}
\textsc{R-Inst} refines only the true branch; \texttt{else} branches keep
the unnarrowed type. Sound (identity), incomplete relative to tsc.

\subsection{Persisting incompleteness: multi-literal \texttt{in} keys}
When the key expression's type is a union of literals, \textsc{R-In} keeps
object members unchanged (still dropping primitives on the true branch);
per-key case splits are not performed.

\subsection{Known unsoundness: string-kind falsiness}
\label{sec:ex-strfalsy}
The lexeme-level falsiness test classifies a string literal by its stored
lexeme (\texttt{lexeme == ""}), but the lexeme is not decoded. An
ECMA-262 \emph{LineContinuation} — a backslash followed by a line
terminator — has string value \texttt{""}, so a non-empty
backslash-newline lexeme denotes the falsy value \texttt{""} while the
test reads it truthy (assumption \textbf{L-LIT}(iii)). This is
\emph{observable}. Let \texttt{L} be such a line-continuation string
literal and take \texttt{x: L | 0} — both members falsy at runtime.
\texttt{L} is classified truthy, so the union appears to have exactly one
falsy member; the boolean-condition guard passes, the falsy residue of
\texttt{if (x)} drops \texttt{L}, and in the else-branch
\texttt{const z: 0 = x} is \emph{accepted} while $x$ holds the string
\texttt{""} at runtime. The canonical control \texttt{"" | 0} is
correctly rejected as a two-falsy-value condition. The sibling assumption
\textbf{L-LIT}(ii) fails on the same undecoded-lexeme root: two spellings
of one string (e.g.\ \texttt{"a"} and \texttt{"\textbackslash x61"}) are
stored as distinct literals, so an equality that holds at runtime is
typed as always-false.

%==========================================================================
\section{Assumptions and Proof Obligations}
\label{sec:assumptions}

Every axiom this document assumes, with its implementation anchor and a
status from a fixed vocabulary:
\textsc{proved} (established in this document or by inspection of the
anchored code), \textsc{assumed} (believed true, undischarged),
\textsc{open} (a stated obligation whose discharge is planned), and
\textsc{violated} (a known counterexample exists in the implementation).
A status may carry a note. Items whose discharge should be a pinned test
say so in the note; wiring those tests is part of the maintenance plan
(Appendix~\ref{app:maintenance}). The obligations of
Section~\ref{sec:inference} (Obligations~\ref{obl:sub}--\ref{obl:restore})
belong to this section's inventory; their current statuses (after the
deep-dive series, \texttt{deep-dives/README.md}):
Obligation~\ref{obl:sub} partially \textsc{proved} with named residues
(deep dive D), Obligation~\ref{obl:monotone} \textsc{violated} in two
witnessed instances (see its status note), and
Obligation~\ref{obl:restore} \textsc{proved} (deep dive E, with a
documented carve-out).

% GENERATED by docs/tools/docs.py from docs/ledger.yaml — do not edit.

\begingroup
\renewcommand{\arraystretch}{1.2}
\small
\begin{longtable}{@{}>{\raggedright\arraybackslash}p{2.5cm}>{\raggedright\arraybackslash}p{8.6cm}>{\raggedright\arraybackslash}p{3.1cm}@{}}
\toprule
\textbf{Item} & \textbf{Statement} & \textbf{Status} \\
\midrule
\endfirsthead
\toprule
\textbf{Item} & \textbf{Statement} & \textbf{Status} \\
\midrule
\endhead
\midrule
\multicolumn{3}{r@{}}{\small\emph{Continued on next page}} \\
\endfoot
\bottomrule
\endlastfoot
L-TAG & Every runtime value carries exactly one tag consistent with its static tag. & \textsc{proved} \\
L-UNION & \texttt{union} over-approximates set union (all levels). & \textsc{proved} \\
L-INTER & When \texttt{inter} returns a non-$\bot$ result, that result contains the semantic intersection of its operands, provided every recursive call that returned $\bot$ had disjoint operands (L-INTER-BOT). & \textsc{open} \\
L-INTER-BOT & A $\bot$ result from \texttt{inter} implies semantic disjointness. & \textsc{violated} \\
L-EXP & Expansion preserves or enlarges semantics (modes \texttt{None}/Covariant); residual instances keep theirs. & \textsc{proved} \\
L-MONO & Legal inference-state changes during refinement preserve membership in types computed earlier. & \textsc{proved} \\
L-SEQ & Type denotations used as equality RHSs are $\approx$-closed. & \textsc{violated} \\
L-ASSERT & \texttt{console.assert}$(e)$ establishes $e$ on the subsequent path. & \textsc{violated} \\
L-SUB & \texttt{Clauses::isSub} acceptance implies semantic inclusion under the current state. & \textsc{open} \\
A-CTOR-CALL & Class-constructor values carry $\mathsf{CPresent}$ construct and $\mathsf{CFunction}$ call modes, consistent with $[[\mathrm{Call}]]$/$[[\mathrm{Construct}]]$. & \textsc{proved} \\
L-TC & Checked expressions inhabit their checked types. & \textsc{assumed} \\
L-LIT & Literal spellings: (i) $\delta$ is a function; (ii) injectivity modulo $\approx$ after normalization; (iii) falsiness-exactness of the lexeme test. & \textsc{violated} \\
L-INST & $P_{\mathsf{inst}}(C) \subseteq \sem{T_C}$ for the instance type as instantiated at the narrowing site. & \textsc{proved} \\
L-ISARR & $\mathsf{isArrayV}(o)$ implies $o \in \sem{\mathsf{ReadonlyArray}\langle\top\rangle}$. & \textsc{proved} \\
L-LOCAL & Module-level bindings are never narrowed; function-level locals are, including captured ones. & \textsc{proved} \\
L-PURE & Evaluating a condition does not mutate the narrowed locals. & \textsc{violated} \\
L-CAPT & Assignments to captured annotated or initializer-inferred locals are checked against the binding's declared or widened inferred type. & \textsc{proved} \\
L-CAPT-EVOLVE & A captured write to an evolving binding (bare \texttt{let x = null/undefined}) is unchecked, and later reads keep the stale evolved type. & \textsc{violated} \\
L-STALE-READ & A closure body invoked later reads enclosing narrowings that are stale at invocation time. & \textsc{violated} \\
L-LOOP-CARRIED & A pre-loop narrowing of an annotated binding survives a back-edge write that invalidates it. & \textsc{violated} \\
L-CLONE & A type computed by \texttt{inter} under a cloned substitution over-approximates the intersection under the main one. & \textsc{proved} \\
L-MUT & Field writes are governed by $\mathsf{set}$-types; get-only membership is sound for snapshot reads. & \textsc{assumed} \\
L-MUT-ALIAS & A checked field or index write through one mutable view preserves the membership guarantees of every live alias of the same runtime object. & \textsc{violated} \\
L-GETSET & Get/set coherence of accessor fields. & \textsc{open} \\
\end{longtable}
\endgroup
\paragraph{Details.}\leavevmode
\medskip\par\noindent\textbf{L-TAG}\quad(\textsc{proved})\\*[0.15em]
Every runtime value carries exactly one tag consistent with its static tag (the runtime counterpart of Lemmas A0/A1, which are proved only inside the model)\\*[0.15em]
{\small\emph{Anchor:} ECMAScript language types / \texttt{typeof} (there is no SKJS code generator; checked programs run on a JS engine)}\par
\nopagebreak
\textsc{proved} (deep dive A: realization map into ECMA-262 \S6.1, modulo the environment assumption A-DDA)

\medskip\par\noindent\textbf{L-UNION}\quad(\textsc{proved})\\*[0.15em]
\texttt{union} over-approximates set union (all levels)\\*[0.15em]
{\small\emph{Anchor:} \texttt{TTypeUnion.sk}}\par
\nopagebreak
\textsc{proved} (deep dive B — in fact exact at every representation layer for the membership semantics)

\medskip\par\noindent\textbf{L-INTER}\quad(\textsc{open})\\*[0.15em]
Non-bottom containment: when $\mathtt{inter}(t_1,t_2) \neq \bot$, $\sem{t_1} \cap \sem{t_2} \subseteq \sem{\mathtt{inter}(t_1,t_2)}$, conditional on L-INTER-BOT for recursive intersections used to build the result. The top-level $\bot$ case is also L-INTER-BOT; a consumer that concludes an inhabitant is \emph{in} the implemented intersection needs both rows, since a $\bot$ result empties the conclusion\\*[0.15em]
{\small\emph{Anchor:} \texttt{TTypeInter.sk}}\par
\nopagebreak
\textsc{proved} (deep dive B) for every fresh structural intersection that returns non-$\bot$, conditional on L-INTER-BOT for recursive calls and using the canonical TVAR-SEM definition for the instance layer; \textsc{open} for stale \texttt{seenInter} reuse under inference growth (C-FREEZE-CACHE)

\medskip\par\noindent\textbf{L-INTER-BOT}\quad(\textsc{violated})\\*[0.15em]
A $\bot$ result from \texttt{inter} implies semantic disjointness — the property required by refinement rules that discard an impossible branch and by L-SUB's covariant variable read (O-GETCOV)\\*[0.15em]
{\small\emph{Anchor:} \texttt{TTypeInter.sk}}\par
\nopagebreak
\textsc{violated} (deep dive B) in two cases: the \texttt{TFun::inter} case — incomparable signatures treated as disjoint (issue report \path{tfun-inter-incomparable-not-disjoint.md}); and the num-kind beyond-precision spelling case — distinct spellings of one double kept distinct, their lexeme intersection a spurious $\bot$ (mediated by L-LIT; issue report \path{beyond-precision-numeric-literals.md}). The per-tag \texttt{STypes} assembly, within-precision numeric spellings, and decoded string spellings produce only semantic bottoms and are \textsc{proved}; stale \texttt{seenInter} reuse under inference growth remains the separate open cache obligation C-FREEZE-CACHE

\medskip\par\noindent\textbf{L-EXP}\quad(\textsc{proved})\\*[0.15em]
Expansion preserves or enlarges semantics; residual instances keep theirs\\*[0.15em]
{\small\emph{Anchor:} \texttt{TType::expand}}\par
\nopagebreak
\textsc{proved} (deep dive C) for the modes refinement uses (\texttt{None}/Covariant) — and shown \emph{false} in contravariant mode (unbound reads produce $\bot$), so the axiom is hereby scoped to those modes; all refinement clients verified to use only them. The containment proof consumes both L-INTER cases and is therefore conditional on L-INTER-BOT (\textsc{violated}) as well as on L-INTER's open cache-reuse case. Reusing an alias expansion is exact because alias normalization does not read existing variable bindings. The remaining conditions are recorded separately as E-NORM, E-ARR, E-CLASS-SYM, E-GUARD, E-EXPKEY, and E-PLACEHOLDER-ATOM; see deep dive C

\medskip\par\noindent\textbf{L-MONO}\quad(\textsc{proved})\\*[0.15em]
Legal inference-state changes during refinement preserve membership in types computed earlier (Axiom~\ref{ax:mono})\\*[0.15em]
{\small\emph{Anchor:} \texttt{Subst::get}, refinement sites}\par
\nopagebreak
\textsc{proved} (deep dive E, given M1 — semantic monotonicity in variable envelopes — and L-UNION); \texttt{setConstraint} is a functional update, so there is no shared-state threat, and the $\mathrm{range}(\Env)$ restriction is droppable (see the axiom's remark)

\medskip\par\noindent\textbf{L-SEQ}\quad(\textsc{violated})\\*[0.15em]
Type denotations used as equality RHSs are $\approx$-closed\\*[0.15em]
{\small\emph{Anchor:} \texttt{evalVarEquality}}\par
\nopagebreak
\textsc{violated}, witnessed end-to-end at the undefined/\texttt{void} corner (\S\ref{sec:ex-voidseq}; issue report \path{void-undefined-equality-never.md}): \texttt{evalVarEquality}'s generic variable-RHS arm installs $\mathtt{inter}(\texttt{void},\texttt{undefined}) = \texttt{never}$ while $\approx$ relates the two units, and the \texttt{void}-operand check on \texttt{===} is a warning, so the corner is reachable — the branch runs with \texttt{undefined} bound at \texttt{never}. The undefined-\emph{literal} arm filters the units jointly and is sound; literal/full primitives and objects remain $\approx$-closed by construction

\medskip\par\noindent\textbf{L-ASSERT}\quad(\textsc{violated})\\*[0.15em]
The premise behind the \texttt{console.assert} refinement clause: the call must not return normally when $e$ is false, and the callee must resolve to the global assertion function — neither condition holds\\*[0.15em]
{\small\emph{Anchor:} \texttt{TypeCheckExpr.sk} (the \texttt{console.assert} match)}\par
\nopagebreak
\textsc{violated}, witnessed end-to-end in both directions (issue draft \path{console-assert-narrowing.md}): a failed \texttt{console.assert} logs and continues, so the narrowed path runs with the condition false; and the match is by spelling alone, so a shadowing binding named \texttt{console} with an arbitrary \texttt{assert} method triggers the same refinement. tsc rejects both witnesses. Fix: remove the refinement or register a recognized assertion function that does not return when its condition is false

\medskip\par\noindent\textbf{L-SUB}\quad(\textsc{open})\\*[0.15em]
\texttt{Clauses::isSub} acceptance implies semantic inclusion under the current state\\*[0.15em]
{\small\emph{Anchor:} \texttt{TTypeSubtyping.sk}}\par
\nopagebreak
split (deep dive D): \textsc{proved} — dispatcher architecture, \texttt{tenv.seenSub} version-keying (corrected argument), coinductive reduction (via SEM-GFP), primitives, \texttt{STypes}, DNF covering, object skeleton, and field reads (conditional on L-SUB-IDXNAME at their numeric-named-field/dual-index corner); conditional — field \emph{set} direction on L-MUT, functions on the E-SIG signature semantics; \textsc{violated} — \texttt{subst.seenSub} results recorded while candidate growth is suppressed (witnessed; issue report \path{freeze-window-stale-memo-caches.md}), and the two brands used by refinement but never validated by subtyping: a callable can enter a type marked \texttt{CAbsent} (F-CABSENT), and the never-compared array brand (F-BRAND)

\medskip\par\noindent\textbf{A-CTOR-CALL}\quad(\textsc{proved})\\*[0.15em]
Class-constructor values are typed with construct mode $\mathsf{CPresent}$ and call mode $\mathsf{CFunction}$ (Definition Callability), so \texttt{typeof}-\texttt{"function"} refinement keeps them — the runtime facts $[[\mathrm{Call}]]$ and $[[\mathrm{Construct}]]$ are both represented\\*[0.15em]
{\small\emph{Anchor:} \texttt{newStaticObject}; \path{narrowByTypeofFunction}}\par
\nopagebreak
\textsc{proved} (deep dive A: constructor values built with both a construct signature and known callability, and the constructible-implies-$[[\mathrm{Call}]]$ invariant enforced at \texttt{TFields} construction); the remaining validation gap for the brand-like $\mathsf{CAbsent}$ claim is the separate row F-CABSENT

\medskip\par\noindent\textbf{L-TC}\quad(\textsc{assumed})\\*[0.15em]
Checked expressions inhabit their checked types (used for the RHS of \textsc{R-Eq} and field reads)\\*[0.15em]
{\small\emph{Anchor:} \texttt{TypeCheckExpr.sk}}\par
\nopagebreak
\textsc{assumed} (axiom here; a full development would discharge it by mutual induction with a preservation theorem, of which the branch-soundness lemmas are the condition case)

\medskip\par\noindent\textbf{L-LIT}\quad(\textsc{violated})\\*[0.15em]
Three claims about literal spellings: (i) $\delta$ is a function — definitional, used by R2; (ii) \emph{injectivity modulo $\approx$} after normalization — needed only for \emph{exactness} claims; (iii) \emph{falsiness-exactness}: the lexeme test \texttt{isFalsy} agrees with Definition~\ref{def:truthiness} — needed by the residue lemma\\*[0.15em]
{\small\emph{Anchor:} \texttt{normalizeJSNumber}/\texttt{normalizeJSBigInt}; \texttt{isFalsy}; string-escape decoding; enum fold outputs}\par
\nopagebreak
(i) \textsc{proved} (definitional). For the \emph{string} kind, (ii)/(iii) are \textsc{proved}: escape sequences are decoded at parse time (\texttt{JSStringEscape.sk}), so a literal's stored lexeme is its denoted string — one spelling per value (ii) and the lexeme emptiness test agrees with falsiness (iii). For num/bigint, (ii)/(iii) are \textsc{proved} for \emph{within-precision} literals: normalized to canonical form at parse time (\texttt{normalizeJSNumber}/\texttt{normalizeJSBigInt}), so on that range zero has one spelling per kind and injectivity holds. For \emph{beyond-precision} numeric literals — those failing the round-trip $\mathtt{String(Number(s))} = s$: 16--17-digit decimals, radix forms $> 2^{53}$, over/underflow exponent forms — (ii)/(iii) are \textsc{violated}: distinct spellings of one double are kept distinct, so \texttt{===} narrowing yields a spurious \texttt{never} (ii), and the normalizer can emit a zero-denoting lexeme (\texttt{1e-400}, underflow to $+0$) classified truthy (iii) — both witnessed end-to-end (issue report \path{beyond-precision-numeric-literals.md}). All three claims quantify over \emph{all stored literals}, enum-folder-synthesized ones included (deep dive G proves the fold outputs canonical in form); whether a synthesized literal equals the runtime value of its source expression is the orthogonal row L-LIT-FOLD. All other results are unconditional

\medskip\par\noindent\textbf{L-INST}\quad(\textsc{proved})\\*[0.15em]
$P_{\mathsf{inst}}(C) \subseteq \sem{T_C}$ for the instance type as instantiated at the narrowing site\\*[0.15em]
{\small\emph{Anchor:} \texttt{evalInstanceof}}\par
\nopagebreak
\textsc{proved} (deep dive A) modulo A-HASINST and A-PROTO (both \textsc{assumed}) and, for intersected branches, L-INTER-BOT (\textsc{violated}): the containment step of the proof needs the intersection to be non-$\bot$ (L-INTER) — when it is $\bot$ the branch is dropped, which is sound only under the disjointness clause. Narrowing preserves the scrutinee's own type arguments — the instance view is built as \texttt{TFields} from the resolved class with the preservation probe run per branch — so the installed type is inhabited by exactly the values \texttt{instanceof} admits at those arguments (skolem interpretation: SKOLEM-SEM); no fresh unbound candidates are introduced at the narrowing site

\medskip\par\noindent\textbf{L-ISARR}\quad(\textsc{proved})\\*[0.15em]
$\mathsf{isArrayV}(o)$ implies $o \in \sem{\mathsf{ReadonlyArray}\langle\top\rangle}$\\*[0.15em]
{\small\emph{Anchor:} \texttt{evalIsArray}}\par
\nopagebreak
\textsc{proved} (deep dive A: spec-level \texttt{IsArray} analysis + verification that the \texttt{\_\_SKJS\_isArray} brand is stripped into the \texttt{isArray} component, never a required field) modulo A-PROXY

\medskip\par\noindent\textbf{L-LOCAL}\quad(\textsc{proved})\\*[0.15em]
Module-level bindings are never narrowed (they resolve through \texttt{defMap}, not the scope chain, so the primitive's early return fires). Function-level locals \emph{are} narrowed — including locals of enclosing functions reached from inside a closure, since scope lookup walks the whole chain\\*[0.15em]
{\small\emph{Anchor:} \texttt{setLocalFilterMap}; \texttt{TypeCheckScope.sk}}\par
\nopagebreak
\textsc{proved} (by code inspection) for the module-level half; the captured-local half sharpens L-PURE; a pinned test is recommended. Narrowing module-level bindings is a tracked future enhancement (issue \#585, blocked on immutability)

\medskip\par\noindent\textbf{L-PURE}\quad(\textsc{violated})\\*[0.15em]
Evaluating a condition does not mutate the narrowed locals\\*[0.15em]
{\small\emph{Anchor:} —}\par
\nopagebreak
\textsc{violated} (empirically pinned): the checker narrows across a call that assigns a captured local and accepts the resulting unsound program (\texttt{if (x !== null \&\& mut())} with \texttt{mut} setting \texttt{x = null}) — the same unsoundness tsc accepts. Tracked as issues \#597 (side-effecting call between test and use) and \#599 (captured local mutated via a lambda). The one-step scoping of the branch-soundness theorems (\S\ref{sec:scope}) is therefore a real boundary, not caution. Deep dive F specifies the fix: a layered invalidation policy (P1--P4) with a preservation lemma replacing the one-step scoping

\medskip\par\noindent\textbf{L-CAPT}\quad(\textsc{proved})\\*[0.15em]
Assignments to captured annotated or initializer-inferred locals are checked against the binding's declared or widened inferred type (an implicit premise of L-TC for closure bodies); evolving bindings are the separate row L-CAPT-EVOLVE\\*[0.15em]
{\small\emph{Anchor:} \texttt{Scope::set}; \texttt{TypeCheckExpr.sk} (lambda checking)}\par
\nopagebreak
\textsc{proved} (deep dive F): assignments to captured annotated or initializer-inferred locals are routed through the binding's write path and checked against the declared or widened inferred type, in lambda bodies as at top level. The evolving-binding exemption is the separate violated row L-CAPT-EVOLVE; the read-side staleness of enclosing narrowings is the separate row L-STALE-READ

\medskip\par\noindent\textbf{L-CAPT-EVOLVE}\quad(\textsc{violated})\\*[0.15em]
Evolving bindings are exempt from the captured-write check, and reads after the closure runs do not incorporate the written value — the part of the captured-write question that L-CAPT's proved statement excludes\\*[0.15em]
{\small\emph{Anchor:} \texttt{Scope::set} (evolving exemption); \texttt{getSingleVar}}\par
\nopagebreak
\textsc{violated}, witnessed end-to-end (deep dive F; issue report \path{evolving-capture-unchecked.md}): a closure writes \texttt{"s"} into a bare \texttt{let x = null} binding and a later read binds the string at type \texttt{null}; tsc accepts the same program (evolving \texttt{any} parity). Fix: track captured writes in the evolved union or check them against it at the write site

\medskip\par\noindent\textbf{L-STALE-READ}\quad(\textsc{violated})\\*[0.15em]
A closure body is checked using an enclosing narrowing even when enclosing code can invalidate that narrowing before invoking the closure; \#599 covers the opposite write-then-read direction\\*[0.15em]
{\small\emph{Anchor:} \texttt{TypeCheckExpr.sk} (lambda checking); narrowed-scope capture}\par
\nopagebreak
\textsc{violated}, witnessed end-to-end in both forms — a captured narrowed local and a narrowed field read inside a closure (deep dive F; issue report \path{stale-closure-reads-narrowing.md}); the fix is deep dive F's layered invalidation policy, shared with L-PURE

\medskip\par\noindent\textbf{L-LOOP-CARRIED}\quad(\textsc{violated})\\*[0.15em]
Loop-carried narrowing staleness: \texttt{Scope::set}'s annotation arm checks a back-edge write only against the declared type, so a pre-loop narrowing of an annotated binding survives into iteration 2 after a write that falsifies it; un-annotated bindings are protected by the loop-write check\\*[0.15em]
{\small\emph{Anchor:} \texttt{Scope::set} (annotation arm); loop merge}\par
\nopagebreak
\textsc{violated}, witnessed end-to-end (deep dive F; issue report \path{loop-carried-narrowing-stale.md}): the narrowed read binds \texttt{null} at \texttt{string} on iteration 2; tsc rejects. Fix fix: deep dive F's P1 loop-entry invalidation layer

\medskip\par\noindent\textbf{L-CLONE}\quad(\textsc{proved})\\*[0.15em]
A type computed by \texttt{inter} under a cloned substitution over-approximates the intersection when later read under the main one\\*[0.15em]
{\small\emph{Anchor:} \texttt{evalVarEquality}}\par
\nopagebreak
\textsc{proved} (deep dive E, fork-state bridging, modulo the fork-state forms of L-INTER and L-INTER-BOT); CL-1 is discharged by a non-escape audit: the \texttt{instanceof} preservation probe generates variables on a clone, but uses them only in probe-side subtype checks, so none reaches a type installed against the main substitution. The property should be guarded by the recommended CI assertion

\medskip\par\noindent\textbf{L-MUT}\quad(\textsc{assumed})\\*[0.15em]
Field writes are governed by $\mathsf{set}$-types via subtyping; membership semantics (get-only) is sound for the snapshot reads refinement performs\\*[0.15em]
{\small\emph{Anchor:} \texttt{TField}}\par
\nopagebreak
\textsc{assumed}; full mutation model out of scope. The tracked instance is issue \#598: field reassignment inside a refined branch does not re-widen the narrowed field, invalidating \textsc{R-EqField}/\textsc{R-TruthyField} results past the write; deep dive F specifies the re-widening rule (W-Field)

\medskip\par\noindent\textbf{L-MUT-ALIAS}\quad(\textsc{violated})\\*[0.15em]
Cross-view mutation coherence, direct writes: a checked field or index write through one mutable view preserves the membership guarantees of every live alias of the same runtime object; call-mediated mutation is the separate row L-MUT-ALIAS-CALL\\*[0.15em]
{\small\emph{Anchor:} \texttt{evalInstanceof}; indexed writes}\par
\nopagebreak
\textsc{violated} (deep dive H) when an alias is created after field narrowing, for writes through either alias (issue report \path{narrowed-alias-stale-write.md}); types introduced directly by \texttt{instanceof} narrowing are coherent — the \texttt{Array} arm preserves the readonly index and the class arm preserves the scrutinee's type arguments — except when a readonly index covers a named field introduced by the class (V-MINT-ROIDX). Aliasing between declaration-derived types is coherent, conditional on L-SUB-FIELDS-SET

\medskip\par\noindent\textbf{L-GETSET}\quad(\textsc{open})\\*[0.15em]
Get/set coherence of accessor fields — the prerequisite for the sound accessor variant of deep dive F's W-Field re-widening rule\\*[0.15em]
{\small\emph{Anchor:} \texttt{TField} (get/set split)}\par
\nopagebreak
\textsc{open} (deep dive F)



\paragraph{Series-introduced items.}
The deep-dive series introduces further named assumptions and obligations
(interface facts its proofs consume, and residues its audits leave). They
are registry entries like the rows above and carry the same status
vocabulary:

% GENERATED by docs/tools/docs.py from docs/ledger.yaml — do not edit.

\begingroup
\renewcommand{\arraystretch}{1.2}
\small
\begin{longtable}{@{}>{\raggedright\arraybackslash}p{2.5cm}>{\raggedright\arraybackslash}p{8.6cm}>{\raggedright\arraybackslash}p{3.1cm}@{}}
\toprule
\textbf{Item} & \textbf{Statement} & \textbf{Status} \\
\midrule
\endfirsthead
\toprule
\textbf{Item} & \textbf{Statement} & \textbf{Status} \\
\midrule
\endhead
\midrule
\multicolumn{3}{r@{}}{\small\emph{Continued on next page}} \\
\endfoot
\bottomrule
\endlastfoot
L-MUT-ALIAS-CALL & Call-mediated mutation (setter-shaped methods, \texttt{push}) through one view preserves the membership guarantees of every live alias. & \textsc{open} \\
V-INTER-SET & Reachability of an incoherent live pair through the intersection's permissive $\mathsf{set}$-union, independent of the already violated admission and minting sites. & \textsc{open} \\
OBL-MONOTONE & Monotone Reads Under Growth: judgements validated under $\rho$ survive later legal transitions. & \textsc{violated} \\
OBL-RESTORE & No Commitment on Failure: failing and exploratory paths restore all relevant state. & \textsc{proved} \\
MR-1 & No \texttt{subst.seenSub} memo recorded under \texttt{freezeCandidates} may outlive growth of an open candidate readable from its left side. & \textsc{violated} \\
C-FREEZE-CACHE & Isolated attribution of \texttt{seenInter} freeze-window staleness (mechanism active; shared witness with MR-1). & \textsc{open} \\
CL-1 & No fresh variable generated on a clone occurs in a type installed against the main substitution. & \textsc{proved} \\
M1 & Semantic monotonicity of $\sem{\cdot}$ in variable envelopes (a property of the get-only, signature-opaque semantics). & \textsc{assumed} \\
I-F & Fork-state form of L-INTER: the containments hold from a forked substitution state (consumed by L-CLONE's proof). & \textsc{assumed} \\
RB-1 & Maintenance condition: \texttt{isSub} callers discarding errors are save/restore-wrapped or unreachable by uncommitted array literals. & \textsc{open} \\
E-NORM & Named-type normalization is semantics-preserving (interface consumed by expansion). & \textsc{assumed} \\
E-ARR & Delayed array-literal forms keep their semantics under expansion. & \textsc{assumed} \\
E-CLASS-SYM & The class path of expansion is binding-blind (or its binding reads are harmless). & \textsc{open} \\
E-GUARD & Recursion-guard identity expansions are never cached into long-lived substitutions. & \textsc{open} \\
SEM-GFP & The set-theoretic semantics is the greatest solution of the recursive membership equations. & \textsc{proved} \\
E-SIG & Behavioral call-soundness semantics for function signatures (required for function subtyping to be a soundness statement). & \textsc{assumed} \\
A-HASH & \texttt{TType::compare} equality is trustworthy (XXHash-only fast paths). & \textsc{assumed} \\
A-PROTO & No prototype surgery / Proxy GetPrototypeOf traps reach checked code (enforced by stub omissions). & \textsc{assumed} \\
A-DDA & No [[IsHTMLDDA]] object (\texttt{document.all}) reaches checked code (vacuous on Node). & \textsc{assumed} \\
A-PROXY & No Proxy exotic object reaches checked code. & \textsc{assumed} \\
A-HASINST & No non-default \texttt{Symbol.hasInstance} on narrowed \texttt{instanceof} RHS classes. & \textsc{violated} \\
TVAR-SEM & A type-variable occurrence denotes its substitution envelope intersected with its occurrence constraints; an unbound envelope is empty. & \textsc{proved} \\
SET-SEM & A view authorizes exactly the writes its resolved $\mathsf{set}$-types accept; a readonly field authorizes nothing. & \textsc{proved} \\
L-SUB-FIELDS-SET & \texttt{isSub}'s field arms preserve write capability downward: acceptance implies each target field's $\mathsf{set}$-type is below the source's, and readonly fields acquire no writers. & \textsc{open} \\
O-GETCOV & Soundness of the covariant variable read in subtyping (inherits L-INTER's violated $\bot$-clause). & \textsc{open} \\
O-GETCONTRA & Soundness of the contravariant variable read in subtyping. & \textsc{open} \\
O-ARRAYLIT & Soundness of the array-literal decision procedure inside subtyping. & \textsc{open} \\
F-CALL & The \texttt{isSubOptTFuns} \texttt{None}-mode arm reads \texttt{call{=}None} as unknown callability (vs.\ refinement's definitely-not-callable). & \textsc{violated} \\
F-BRAND & Subtyping never compares the array brand, though refinement narrows on it. & \textsc{violated} \\
SETCONSTRAINT-FUN & \texttt{TVar::setConstraint} is a functional update on a frozen value: constraints are occurrence-level, not shared state. & \textsc{proved} \\
E-TRANSPORT & Monotone-reads transport: an inclusion validated under $\rho$ survives later legal transitions on the same timeline, modulo MR-1. & \textsc{proved} \\
TENV-EXPANDINTER & The \texttt{tenv.expandInter} mode flag is path-scoped, and both intersection forms it selects (expanded and deferred \texttt{unexpandedInter}) satisfy the L-INTER containment. & \textsc{open} \\
L-SUB-IDXNAME & A numeric-named field is checked against both a number and a string index when the supertype has both (soundness of the width lemma's numeric-named case). & \textsc{open} \\
A-INTERSEM & $\sem{\mathsf{TInter}(\bar t)} \subseteq \bigcap \sem{t}$ — the instance-layer upper bound the dispatcher's fast paths assume. & \textsc{assumed} \\
E-GUARD-PLACEHOLDER & An un-annotated ambient-class field placeholder must not head-expand to $\bot$ (the case E-GUARD's recursion-guard audit did not cover). & \textsc{violated} \\
ACCESSOR-METHOD & Accessor (\texttt{get}/\texttt{set}) members must be typed as fields, not methods; today a getter is a call signature, so \texttt{c.p()} type-checks and throws. & \textsc{violated} \\
L-SUB-TSINST & The reflexive \texttt{TSInst} fast path in \texttt{isSub} (same name, arguments equal after \texttt{peekRaw}) is sound. & \textsc{violated} \\
\end{longtable}
\endgroup
\paragraph{Details.}\leavevmode
\medskip\par\noindent\textbf{L-MUT-ALIAS-CALL}\quad(\textsc{open})\\*[0.15em]
Cross-view mutation coherence, call-mediated facet: a checked method call through one view preserves every live alias's membership; needs a receiver-effect strengthening of E-SIG, which as stated concerns behavioral signature semantics only\\*[0.15em]
{\small\emph{Anchor:} method calls on narrowed receivers}\par
\nopagebreak
\textsc{open} (deep dive H): the minted \texttt{Array}$\langle\top\rangle$ view exposes the same incoherence through \texttt{push} as through the graded index write, but no signature semantics with receiver-heap effects exists to grade it against

\medskip\par\noindent\textbf{V-INTER-SET}\quad(\textsc{open})\\*[0.15em]
\texttt{TField::inter} unions $\mathsf{set}$-types; the open question is whether an incoherent live pair is reachable through it on a route not already graded by the narrowing-admission or minted-view sites\\*[0.15em]
{\small\emph{Anchor:} \texttt{TField::inter} (set arm)}\par
\nopagebreak
\textsc{open} (deep dive H): the mechanism is established; every declared route probed is blocked by the admission gate (the union widens the target's $\mathsf{set}$, making admission harder), and refinement-installed intersections reduce to the graded sites; ambient-declared intersection types enter through the stub trust boundary

\medskip\par\noindent\textbf{OBL-MONOTONE}\quad(\textsc{violated})\\*[0.15em]
Monotone Reads Under Growth: judgements validated under $\rho$ survive later legal transitions.\\*[0.15em]
{\small\emph{Anchor:} \texttt{Subst::get}; \texttt{TType::isSub}}\par
\nopagebreak
\textsc{violated} in two witnessed instances — the freeze-window memo staleness MR-1 and the unsealed \texttt{peekRaw} read of the reflexive \texttt{TSInst} arm (L-SUB-TSINST, not memo-mediated); the general transport statement is the separate \textsc{proved} row E-TRANSPORT

\medskip\par\noindent\textbf{OBL-RESTORE}\quad(\textsc{proved})\\*[0.15em]
No Commitment on Failure: failing and exploratory paths restore all relevant state.\\*[0.15em]
{\small\emph{Anchor:} \texttt{SavedSubst}; save/restore sites}\par
\nopagebreak
\textsc{proved} (deep dive E: 14-row save/restore audit; \texttt{SavedSubst} completeness; array-literal carve-out sound under RB-1)

\medskip\par\noindent\textbf{MR-1}\quad(\textsc{violated})\\*[0.15em]
No \texttt{subst.seenSub} memo recorded under \texttt{freezeCandidates} may outlive growth of an open candidate readable from its left side.\\*[0.15em]
{\small\emph{Anchor:} \texttt{Subst::get} (freeze arm); \texttt{TType::isSub} memo}\par
\nopagebreak
\textsc{violated}, witnessed end-to-end (order-dependent acceptance; deep dive E; issue draft \path{freeze-window-stale-memo-caches.md})

\medskip\par\noindent\textbf{C-FREEZE-CACHE}\quad(\textsc{open})\\*[0.15em]
Isolated attribution of \texttt{seenInter} freeze-window staleness (mechanism active; shared witness with MR-1).\\*[0.15em]
{\small\emph{Anchor:} \texttt{seenInter}; \texttt{TType::inter}}\par
\nopagebreak
\textsc{open} for the isolated statement; mechanism demonstrated active (deep dive B); the MR-1 fix covers both memos

\medskip\par\noindent\textbf{CL-1}\quad(\textsc{proved})\\*[0.15em]
No fresh variable generated on a clone occurs in a type installed against the main substitution.\\*[0.15em]
{\small\emph{Anchor:} \texttt{Subst::clone}, \texttt{Subst::gen}}\par
\nopagebreak
discharged (deep dive E): unreachable — no \texttt{gen*} call site in the clone-compute-install tree; guard with the recommended CI assertion

\medskip\par\noindent\textbf{M1}\quad(\textsc{assumed})\\*[0.15em]
Semantic monotonicity of $\sem{\cdot}$ in variable envelopes (a property of the get-only, signature-opaque semantics).\\*[0.15em]
{\small\emph{Anchor:} main paper semantics}\par
\nopagebreak
\textsc{assumed} (deep dive E; provable by structural induction in the main paper)

\medskip\par\noindent\textbf{I-F}\quad(\textsc{assumed})\\*[0.15em]
Fork-state form of L-INTER: the containments hold from a forked substitution state (consumed by L-CLONE's proof).\\*[0.15em]
{\small\emph{Anchor:} \texttt{TTypeInter.sk}}\par
\nopagebreak
\textsc{assumed} (deep dives B and E); folds into L-INTER's discharge

\medskip\par\noindent\textbf{RB-1}\quad(\textsc{open})\\*[0.15em]
Maintenance condition: \texttt{isSub} callers discarding errors are save/restore-wrapped or unreachable by uncommitted array literals.\\*[0.15em]
{\small\emph{Anchor:} \texttt{ArrayLitTyping}; \texttt{isSub} probe sites}\par
\nopagebreak
\textsc{open} maintenance condition (deep dive E); CI grep recommended

\medskip\par\noindent\textbf{E-NORM}\quad(\textsc{assumed})\\*[0.15em]
Named-type normalization is semantics-preserving (interface consumed by expansion).\\*[0.15em]
{\small\emph{Anchor:} \texttt{TEnv::normalize}}\par
\nopagebreak
\textsc{assumed} (deep dive C)

\medskip\par\noindent\textbf{E-ARR}\quad(\textsc{assumed})\\*[0.15em]
Delayed array-literal forms keep their semantics under expansion.\\*[0.15em]
{\small\emph{Anchor:} \texttt{TArrayLit}}\par
\nopagebreak
\textsc{assumed} (deep dive C)

\medskip\par\noindent\textbf{E-CLASS-SYM}\quad(\textsc{open})\\*[0.15em]
The class path of expansion is binding-blind (or its binding reads are harmless).\\*[0.15em]
{\small\emph{Anchor:} \texttt{newObject}/\texttt{addFields}}\par
\nopagebreak
\textsc{open} (deep dive C; proof and counterexample both outstanding)

\medskip\par\noindent\textbf{E-GUARD}\quad(\textsc{open})\\*[0.15em]
Recursion-guard identity expansions are never cached into long-lived substitutions.\\*[0.15em]
{\small\emph{Anchor:} \texttt{seenExpansion}; recursion guards}\par
\nopagebreak
\textsc{open} (deep dive C); cache-write invariant + pinned test recommended

\medskip\par\noindent\textbf{SEM-GFP}\quad(\textsc{proved})\\*[0.15em]
Greatest-fixpoint reading: $\sem{\cdot}$ is the greatest solution of the recursive membership equations.\\*[0.15em]
{\small\emph{Anchor:} main paper object semantics}\par
\nopagebreak
\textsc{proved} by definition in the main paper (every recursive occurrence in the membership conditions is positive, so the greatest solution exists); deep dive D supplies the stratified characterization used by the coinductive reduction

\medskip\par\noindent\textbf{E-SIG}\quad(\textsc{assumed})\\*[0.15em]
Behavioral call-soundness semantics for function signatures (required for function subtyping to be a soundness statement).\\*[0.15em]
{\small\emph{Anchor:} main paper object semantics (extension)}\par
\nopagebreak
\textsc{assumed} (deep dive D; proposed extension of the main paper's object semantics)

\medskip\par\noindent\textbf{A-HASH}\quad(\textsc{assumed})\\*[0.15em]
\texttt{TType::compare} equality is trustworthy (XXHash-only fast paths).\\*[0.15em]
{\small\emph{Anchor:} \texttt{TType} hashing}\par
\nopagebreak
\textsc{assumed} (deep dive D)

\medskip\par\noindent\textbf{A-PROTO}\quad(\textsc{assumed})\\*[0.15em]
No prototype surgery / Proxy GetPrototypeOf traps reach checked code (enforced by stub omissions).\\*[0.15em]
{\small\emph{Anchor:} import-stubs}\par
\nopagebreak
\textsc{assumed} (deep dive A; enforced for checked code by stub omissions)

\medskip\par\noindent\textbf{A-DDA}\quad(\textsc{assumed})\\*[0.15em]
No [[IsHTMLDDA]] object (\texttt{document.all}) reaches checked code (vacuous on Node).\\*[0.15em]
{\small\emph{Anchor:} environment}\par
\nopagebreak
\textsc{assumed} (deep dive A)

\medskip\par\noindent\textbf{A-PROXY}\quad(\textsc{assumed})\\*[0.15em]
No Proxy exotic object reaches checked code.\\*[0.15em]
{\small\emph{Anchor:} environment}\par
\nopagebreak
\textsc{assumed} (deep dive A)

\medskip\par\noindent\textbf{A-HASINST}\quad(\textsc{violated})\\*[0.15em]
No non-default \texttt{Symbol.hasInstance} on narrowed \texttt{instanceof} RHS classes.\\*[0.15em]
{\small\emph{Anchor:} \texttt{isInstanceofRhs}}\par
\nopagebreak
\textsc{violated}, tsc-parity (deep dive A; issue draft \path{symbol-hasinstance-narrowing.md})

\medskip\par\noindent\textbf{TVAR-SEM}\quad(\textsc{proved})\\*[0.15em]
Canonical \texttt{TVar} denotation: an occurrence denotes its substitution envelope intersected with its occurrence constraints; an unbound envelope is empty.\\*[0.15em]
{\small\emph{Anchor:} \texttt{Subst::get}/\texttt{peek}}\par
\nopagebreak
\textsc{proved} by definition in the main paper and matched to \texttt{Subst::get}: missing bindings read as bottom contravariantly and seal to top before a covariant expansion; occurrence constraints are intersected after either read

\medskip\par\noindent\textbf{SET-SEM}\quad(\textsc{proved})\\*[0.15em]
Write authorization: a view authorizes the write of $v$ at key $p$ iff $v$ inhabits the resolved $\mathsf{set}$-type of $p$ (intersected named-field $\mathsf{set}$-types, else the covering index's, resolved as at lvalue sites); readonly ($\mathsf{set}=\bot$) and keys with no writable field or index authorize nothing; call-mediated writes are E-SIG's\\*[0.15em]
{\small\emph{Anchor:} \path{TFields::resolveFieldType} (\texttt{isLValue})}\par
\nopagebreak
\textsc{proved} by definition in the main paper and matched to the checker's write path: assignment checks the right-hand side against the resolved $\mathsf{set}$-type by subtyping

\medskip\par\noindent\textbf{L-SUB-FIELDS-SET}\quad(\textsc{open})\\*[0.15em]
Write-capability preservation: field acceptance requires $f_2.\mathsf{set} \leq f_1.\mathsf{set}$ (contravariant) and the readonly gate, so under SET-SEM an admitted view authorizes only writes the source view authorizes\\*[0.15em]
{\small\emph{Anchor:} \texttt{TFields::isSub\_}}\par
\nopagebreak
\textsc{open} (conditional): the syntactic checks are verified at every field arm (deep dive D, write-direction lemma); the SET-SEM containment they induce is conditional on the recursive \texttt{isSub} verdicts, exactly as the read direction is inside L-SUB

\medskip\par\noindent\textbf{O-GETCOV}\quad(\textsc{open})\\*[0.15em]
Soundness of the covariant variable read in subtyping (inherits L-INTER's violated $\bot$-clause).\\*[0.15em]
{\small\emph{Anchor:} \texttt{Subst::get} (covariant); \texttt{expandToClauses}}\par
\nopagebreak
\textsc{open} (deep dive D): the covariant variable read applies \texttt{Subst::get}'s constraint-intersection loop, whose soundness rests on the \texttt{L-INTER-BOT} bottoms; its live strata (incomparable signatures; string-spelling; beyond-precision num-spelling) leave the read unproven. On L-SUB's critical path

\medskip\par\noindent\textbf{O-GETCONTRA}\quad(\textsc{open})\\*[0.15em]
Soundness of the contravariant variable read in subtyping.\\*[0.15em]
{\small\emph{Anchor:} \texttt{Subst::get} (contravariant)}\par
\nopagebreak
\textsc{open} (deep dive D)

\medskip\par\noindent\textbf{O-ARRAYLIT}\quad(\textsc{open})\\*[0.15em]
Soundness of the array-literal decision procedure inside subtyping.\\*[0.15em]
{\small\emph{Anchor:} \path{ArrayLitTyping::subtyping}}\par
\nopagebreak
\textsc{open} (deep dive D; delegated)

\medskip\par\noindent\textbf{F-CALL}\quad(\textsc{violated})\\*[0.15em]
The \texttt{isSubOptTFuns} \texttt{None}-mode arm reads \texttt{call{=}None} as unknown callability (vs.\ refinement's definitely-not-callable).\\*[0.15em]
{\small\emph{Anchor:} \texttt{isSubOptTFuns}}\par
\nopagebreak
\textsc{violated}, witnessed end-to-end (deep dive D): a callable assigned to \texttt{\{\}} then narrowed by \texttt{typeof} to \texttt{never}; issue draft \path{class-constructor-call-none-typeof.md}

\medskip\par\noindent\textbf{F-BRAND}\quad(\textsc{violated})\\*[0.15em]
Subtyping never compares the array brand, though refinement narrows on it.\\*[0.15em]
{\small\emph{Anchor:} \texttt{TFields::isSub\_}}\par
\nopagebreak
\textsc{violated} by code-level derivation (deep dive D; composite witness sketch; no end-to-end run)

\medskip\par\noindent\textbf{SETCONSTRAINT-FUN}\quad(\textsc{proved})\\*[0.15em]
\texttt{TVar::setConstraint} is a functional update on a frozen value: constraints are occurrence-level, not shared state.\\*[0.15em]
{\small\emph{Anchor:} \texttt{TVar::setConstraint}}\par
\nopagebreak
\textsc{proved} (deep dive E, by code inspection: frozen class, functional update, both call sites use the return value)

\medskip\par\noindent\textbf{E-TRANSPORT}\quad(\textsc{proved})\\*[0.15em]
Monotone-reads transport: an inclusion validated under $\rho$ survives later legal transitions on the same timeline, modulo MR-1.\\*[0.15em]
{\small\emph{Anchor:} \texttt{Subst} transitions}\par
\nopagebreak
\textsc{proved} (deep dive E, the monotone-reads theorem, given M1 and the MR/BR audit; residue MR-1)

\medskip\par\noindent\textbf{TENV-EXPANDINTER}\quad(\textsc{open})\\*[0.15em]
The \texttt{tenv.expandInter} mode flag is path-scoped, and both intersection forms it selects (expanded and deferred \texttt{unexpandedInter}) satisfy the L-INTER containment.\\*[0.15em]
{\small\emph{Anchor:} \texttt{TEnv.expandInter}; \texttt{TInter::expand}}\par
\nopagebreak
\textsc{open} — found by external review: the flag is absent from deep dive E's unsaved-state inventory; path-scoping should follow the frozen-TEnv threading argument, and the deferred form needs deep dive B's containment coverage

\medskip\par\noindent\textbf{L-SUB-IDXNAME}\quad(\textsc{open})\\*[0.15em]
A numeric-named field is checked against both a number and a string index when the supertype has both (soundness of the width lemma's numeric-named case).\\*[0.15em]
{\small\emph{Anchor:} \texttt{indexForFieldName}; \texttt{TFields::isSub\_}}\par
\nopagebreak
\textsc{open} — found by external review: the width lemma checks a numeric-named sub field against only one super index (\texttt{indexForFieldName} returns a single index), while covering semantics needs both; SKJS enforces neither TS2411 nor TS2413. The subtyping over-acceptance is real but \emph{not} end-to-end unsound at the pin — reads of numeric-string keys stay conservative (\texttt{string|number}), masking it, as verified against the binary; a validator-completeness gap, not a witnessed unsoundness. Cited at deep dive D's width lemma as a named residual

\medskip\par\noindent\textbf{A-INTERSEM}\quad(\textsc{assumed})\\*[0.15em]
$\sem{\mathsf{TInter}(\bar t)} \subseteq \bigcap \sem{t}$ — the instance-layer upper bound the dispatcher's fast paths assume.\\*[0.15em]
{\small\emph{Anchor:} \texttt{TInter::expand}}\par
\nopagebreak
\textsc{assumed} (deep dive D; the instance-layer counterpart of deep dive C's \texttt{TInter} expansion — an upper bound on the denotation, distinct from L-EXP's over-approximation direction)

\medskip\par\noindent\textbf{E-GUARD-PLACEHOLDER}\quad(\textsc{violated})\\*[0.15em]
An un-annotated ambient-class field placeholder must not head-expand to $\bot$ (the case E-GUARD's recursion-guard audit did not cover).\\*[0.15em]
{\small\emph{Anchor:} \texttt{instantiate} (missing-annotation return); \texttt{expandToVectors}}\par
\nopagebreak
\textsc{violated}, witnessed end-to-end (deep dive C): the missing-annotation identity return fires with \texttt{seenInstantiate} empty, the identity is dropped by the expansion seen-set, and the field reads as $\bot$ (accepted into \texttt{never}); issue draft \path{ambient-unannotated-field-bottom.md}. Sibling of the recursion-guard E-GUARD row, which is empirically safe

\medskip\par\noindent\textbf{ACCESSOR-METHOD}\quad(\textsc{violated})\\*[0.15em]
Accessor (\texttt{get}/\texttt{set}) members must be typed as fields, not methods; today a getter is a call signature, so \texttt{c.p()} type-checks and throws.\\*[0.15em]
{\small\emph{Anchor:} \texttt{isGetSet} (parser-only); member typing}\par
\nopagebreak
\textsc{violated}, witnessed end-to-end (SKJS-specific; tsc rejects with TS6234): the parser's \texttt{isGetSet} marker is never read by the checker, so a getter \texttt{get p(): T} is typed as a method \texttt{p(): T}. \texttt{c.p()} is accepted and throws at runtime (\texttt{c.p} is not a function). Issue draft \path{getter-typed-as-method.md}. Adjacent to F's accessor caveat and L-GETSET, but a distinct (member-typing) hole

\medskip\par\noindent\textbf{L-SUB-TSINST}\quad(\textsc{violated})\\*[0.15em]
The reflexive \texttt{TSInst} fast path in \texttt{isSub} (same name, arguments equal after \texttt{peekRaw}) is sound.\\*[0.15em]
{\small\emph{Anchor:} \texttt{TType::isSub} (TSInst reflexive arm)}\par
\nopagebreak
\textsc{violated} (deep dive D): the inclusion at the verdict state holds for unconstrained bindings, but the arm reads an unsealed \texttt{SVOpen} candidate via \texttt{peekRaw} without the covariant seal the slow path performs, so a later contravariant grow widens the binding and reads through aliases break — witnessed end-to-end (a value typed \texttt{number} holds a string; issue draft \path{reflexive-tsinst-fastpath-unsealed-read.md}). For constrained variables the equality claim also fails as stated: the occurrence denotes $\sem{t_2}_\rho \cap \bigcap\sem{C}$, not $\sem{t_2}_\rho$



In summary, after the deep-dive series and its adversarial verification
(series map in \texttt{deep-dives/README.md}; the authoritative per-item
statuses are the tables above, from which this grouping is generated):

% GENERATED by docs/tools/docs.py from docs/ledger.yaml — do not edit.

\begin{itemize}
\item \emph{Proved}: L-TAG; L-UNION; L-INTER's containment (all strata); L-EXP (mode-scoped); L-MONO; the majority strata of L-SUB; A-CTOR-CALL; L-INST (modulo A-HASINST, A-PROTO); L-ISARR; L-CAPT; L-CLONE; OBL-RESTORE; SEM-GFP; TVAR-SEM; SET-SEM.
\item \emph{Assumed}: L-TC (the capstone); L-MUT.
\item \emph{Open}: L-SUB's conditional strata; L-GETSET; the isolated \texttt{seenInter} attribution (C-FREEZE-CACHE); E-CLASS-SYM; E-GUARD; O-GETCOV, O-GETCONTRA, O-ARRAYLIT; TENV-EXPANDINTER; L-SUB-IDXNAME; the ambient-field placeholder atom (E-PLACEHOLDER-ATOM).
\item \emph{Violated, with an executed witness} (each with a tracked issue or an issue draft under \texttt{deep-dives/issues/}): the $\bot{\Rightarrow}$disjoint clause L-INTER-BOT (\texttt{TFun} + num spellings); L-SEQ's undefined/\texttt{void} corner; L-LIT(ii)/(iii) num beyond-precision spellings; the enum folder's concat-minted literals (L-LIT-FOLD); L-PURE (\#597/\#599) with the L-MUT instance \#598; stale closure reads of enclosing narrowings (L-STALE-READ); loop-carried narrowing staleness (L-LOOP-CARRIED); cross-view mutation coherence L-MUT-ALIAS (narrowing-admitted aliases); the two witnessed instances of OBL-MONOTONE (MR-1; the reflexive-\texttt{TSInst} read); the laundered \texttt{CAbsent} brand (F-CABSENT); the expansion-cache key collision (E-EXPKEY); accessor members typed as methods (ACCESSOR-METHOD); L-SUB-TSINST (unsealed \texttt{peekRaw} read in the reflexive \texttt{TSInst} arm).
\item \emph{Violated by code-level derivation} (no end-to-end run yet): the never-compared array brand (F-BRAND).
\end{itemize}


\noindent The violations remain local to leaf
assumptions: every fix restores its assumption without changing a theorem
statement or proof in this document. Two design decisions gate several
fixes and are recorded as fundamental holes rather than bugs: the
\texttt{call{=}None} reading dichotomy and the refinement-invalidation
policy (deep dive F). Known unsoundness here is enumerated, witnessed,
and quarantined — future work, not architectural debt.

\medskip
% GENERATED by docs/tools/docs.py from docs/ledger.yaml — do not edit.

\noindent\emph{None recorded.}


%==========================================================================
\appendix

\section{Implementation Correspondence}
\label{app:impl}

Pinned to \texttt{\skjspinhash}.

\begin{center}
\renewcommand{\arraystretch}{1.25}
\begin{tabular}{@{}l@{\quad}l@{}}
\toprule
\textbf{Model} & \textbf{Implementation} \\
\midrule
$\Tags$ (9 tags) & \texttt{STypeTag} (TType.sk) \\
$\mathsf{unit}, \mathsf{prim}, \mathsf{object}$ &
  \texttt{TPrimUnit}, \texttt{TPrimValue}, \texttt{TObject} (TType.sk) \\
literal payloads & \texttt{PrimLiterals} \\
object branches $B$, fields & \texttt{ArraySet<TFields>}, \texttt{TField} \\
sig modes $\mathsf{None}/(\mathit{opt},\sigma)$ &
  \texttt{TFields.call/new : ?(Bool, TFuns)} \\
array brand & \texttt{TFields.isArray : ?Bool} \\
phantom fields & \texttt{FAPhantom} \\
$\STypes$ & \texttt{STypes} (tag-keyed map) \\
$\clause{V}{S}$, $\Clauses$, $\TType$ &
  \texttt{Clause}, \texttt{Clauses}, \texttt{TType} \\
$\refineop^{g,d}$ & \texttt{setLocalFilterMap} + \texttt{TType::mapSTypes} \\
master list dispatch & \texttt{evalCond} (TypeCheckRefine.sk) \\
$\mathsf{rmc}$ & \texttt{TType::removeConstant} \\
$\mathsf{tres}_b, \mathsf{vres}_b$ &
  \texttt{truthinessResidue}, \texttt{truthinessValueResidue} \\
scope merge $\sqcup$ & \texttt{TCEnv::mergeScope} \\
$\rho$, transitions & \texttt{Subst}, \texttt{SVOpen}/\texttt{SVClosed},
  \texttt{grow}/\texttt{get} (TTypeSubst.sk) \\
$\mathsf{freeze}$ & \texttt{Subst.freezeCandidates} \\
growth pre-pass & \texttt{growTVars} (TTypeSubtyping.sk) \\
validation & \texttt{Clauses::isSub} \\
\bottomrule
\end{tabular}
\end{center}

\section{Maintenance}
\label{app:maintenance}

An advisory maintenance rule (``when the implementation adds a new
refinement, add it to this list'') does not survive contact with
development. The replacement is enforced coupling; the parts marked
\emph{in place} are implemented by \texttt{docs/tools/docs.py}
(\texttt{gen} / \texttt{check} / \texttt{build}):

\begin{enumerate}
\item \textbf{The registry} (\emph{in place}). Statuses live only in
  \texttt{docs/ledger.yaml}; the tables of
  Section~\ref{sec:assumptions}, their summary, and the deep dives'
  rows tables are generated from it (\texttt{docs/generated/}). A status
  change is a one-line registry edit; no document states another
  document's status.
\item \textbf{Consistency checks} (\emph{in place}). The \texttt{check}
  command verifies: registry integrity (status vocabulary, anchor files
  exist, issue references resolve); every issue draft is referenced;
  no temporal/history markers in any document; no numeric cross-document
  references (deep dives must cite main-paper items by \emph{name});
  every document's date goes through the shared pin
  (\texttt{docs/shared/pin.tex}).
\item \textbf{Local builds} (\emph{in place}). The \texttt{build}
  command compiles every document (two \texttt{pdflatex} passes), keeps
  each \texttt{.pdf} next to its source, fails on errors or unresolved
  references, and keeps auxiliary files in the git-ignored
  \texttt{docs/.texbuild/} for warm incremental recompilation.
\item \textbf{Anchor cross-check} (\emph{pending}). Each dispatch arm of
  \texttt{evalCond} carries a comment \texttt{// paper: R-\dots}\ naming
  its rule; the check verifies both directions (every arm has a rule;
  every master-list rule has a live arm). Existence, not correctness —
  semantic drift within an arm still requires review, but silent
  additions and removals become failures.
\item \textbf{Witness repros} (\emph{in place for issue witnesses}).
  Each executed unsoundness witness is a runnable
  \texttt{deep-dives/issues/repros/*.ts} with a header verdict;
  \texttt{docs.py repro} runs each against the pinned checker and asserts
  it, and \texttt{check} enforces that every repro is referenced by an
  issue draft and vice versa. A repro that flips verdict when its fix
  lands fails, forcing a documentation update in the same change. Still
  \emph{pending}: expect-error pins for the incompleteness worked
  examples in the test suite proper.
\item \textbf{Pinned revision.} \texttt{docs/shared/pin.tex} names the
  modeled commit once, for every document. A change that trips any check
  updates the pin as part of restoring it.
\end{enumerate}

\noindent The review protocol for this corpus — scopes and checks for
reviewing this paper, the deep dives, and the issue drafts — lives in
\texttt{docs/README.md}.

\end{document}
